% ============================================================================
%  Newton Sketch adaptativo: apuntes de estudio
%  Archivo unico y autocontenido. Compilar dos veces con pdflatex (o latexmk -pdf).
%  Paquetes: babel (spanish), mathpazo, amsmath, tcolorbox, tikz, pgfplots, algorithm,
%  algpseudocode, listings, booktabs, hyperref (todos en TeX Live / MiKTeX estandar).
% ============================================================================
\documentclass[11pt,a4paper,oneside]{book}

% ---------- Codificación, idioma y tipografía ----------
\usepackage[utf8]{inputenc}
\usepackage[T1]{fontenc}
\IfFileExists{spanish.ldf}{%
  \usepackage[spanish,es-tabla,es-nodecimaldot,es-noshorthands,es-noindentfirst]{babel}%
}{% Respaldo para instalaciones de TeX sin babel-spanish: usa los datos de idioma genéricos de babel
  \usepackage{babel}\babelprovide[import,main]{spanish}%
  \setlocalecaption{spanish}{table}{Tabla}\setlocalecaption{spanish}{listtable}{Índice de tablas}%
  \setlocalecaption{spanish}{contents}{Índice}%
}
\usepackage[sc]{mathpazo}
\linespread{1.05}
\usepackage{microtype}
\usepackage[a4paper,margin=2.5cm,headheight=14pt]{geometry}

% ---------- Matemáticas ----------
\usepackage{amsmath,amssymb,amsthm,mathtools,bm}

% ---------- Gráficos, tablas, algoritmos y código ----------
\usepackage{xcolor}
\usepackage{graphicx}
\usepackage{booktabs,array,multirow}
\usepackage{enumitem}
\usepackage{tikz}
\usetikzlibrary{arrows.meta,positioning,calc,shapes.geometric,fit,backgrounds}
\usepackage{pgfplots}
\pgfplotsset{compat=1.16}
\usepgfplotslibrary{groupplots}
\usepackage[most]{tcolorbox}
\usepackage{algorithm}
\usepackage{algpseudocode}
\usepackage{listings}
\usepackage{fancyhdr}
\usepackage[colorlinks=true,linkcolor=blue!50!black,citecolor=green!40!black,urlcolor=blue!60!black]{hyperref}
\hypersetup{pdftitle={Newton Sketch adaptativo: apuntes de estudio},pdfsubject={Sketching aleatorio, autoconcordancia y despacho economico estocastico}}

% ---------- Colores ----------
\definecolor{azul}{RGB}{28,69,135}
\definecolor{verde}{RGB}{30,110,60}
\definecolor{rojo}{RGB}{160,40,40}
\definecolor{ocre}{RGB}{176,112,20}
\definecolor{gris}{RGB}{90,90,90}

% ---------- Entornos tipo teorema ----------
\newtheorem{teorema}{Teorema}[chapter]
\newtheorem{lema}[teorema]{Lema}
\newtheorem{proposicion}[teorema]{Proposición}
\newtheorem{corolario}[teorema]{Corolario}
\theoremstyle{definition}
\newtheorem{definicion}[teorema]{Definición}
\newtheorem{ejemplo}[teorema]{Ejemplo}
\newtheorem{ejercicio}{Ejercicio}[chapter]
\theoremstyle{remark}
\newtheorem{observacion}[teorema]{Observación}

\tcolorboxenvironment{teorema}{enhanced,breakable,colback=azul!4,colframe=azul!70,boxrule=0.5pt,
  left=6pt,right=6pt,top=4pt,bottom=4pt,before skip=10pt,after skip=10pt}
\tcolorboxenvironment{lema}{enhanced,breakable,colback=azul!2,colframe=azul!45,boxrule=0.4pt,
  left=6pt,right=6pt,top=4pt,bottom=4pt,before skip=10pt,after skip=10pt}
\tcolorboxenvironment{proposicion}{enhanced,breakable,colback=azul!2,colframe=azul!45,boxrule=0.4pt,
  left=6pt,right=6pt,top=4pt,bottom=4pt,before skip=10pt,after skip=10pt}
\tcolorboxenvironment{corolario}{enhanced,breakable,colback=azul!2,colframe=azul!45,boxrule=0.4pt,
  left=6pt,right=6pt,top=4pt,bottom=4pt,before skip=10pt,after skip=10pt}
\tcolorboxenvironment{definicion}{enhanced,breakable,colback=verde!4,colframe=verde!60,boxrule=0.4pt,
  left=6pt,right=6pt,top=4pt,bottom=4pt,before skip=10pt,after skip=10pt}

% Cajas de apoyo al estudio
\newtcolorbox{ideaclave}[1][Idea clave]{enhanced,breakable,colback=ocre!7,colframe=ocre!80,boxrule=0.5pt,
  fonttitle=\bfseries\sffamily,title={#1},left=6pt,right=6pt,top=4pt,bottom=4pt,before skip=10pt,after skip=10pt}
\newtcolorbox{resumen}[1][Resumen del capítulo]{enhanced,breakable,colback=gris!6,colframe=gris!70,boxrule=0.5pt,
  fonttitle=\bfseries\sffamily,title={#1},left=6pt,right=6pt,top=4pt,bottom=4pt,before skip=12pt,after skip=10pt}
\newtcolorbox{cuidado}[1][Cuidado]{enhanced,breakable,colback=rojo!5,colframe=rojo!70,boxrule=0.5pt,
  fonttitle=\bfseries\sffamily,title={#1},left=6pt,right=6pt,top=4pt,bottom=4pt,before skip=10pt,after skip=10pt}

% ---------- Algoritmos en español ----------
\floatname{algorithm}{Algoritmo}
\renewcommand{\listalgorithmname}{Lista de algoritmos}
\algrenewcommand\algorithmicrequire{\textbf{Entrada:}}
\algrenewcommand\algorithmicensure{\textbf{Salida:}}
\algrenewcommand\algorithmicfor{\textbf{para}}
\algrenewcommand\algorithmicdo{\textbf{hacer}}
\algrenewcommand\algorithmicwhile{\textbf{mientras}}
\algrenewcommand\algorithmicif{\textbf{si}}
\algrenewcommand\algorithmicthen{\textbf{entonces}}
\algrenewcommand\algorithmicelse{\textbf{si no}}
\algrenewcommand\algorithmicend{\textbf{fin}}
\algrenewcommand\algorithmicreturn{\textbf{devolver}}
\algrenewcommand\algorithmicrepeat{\textbf{repetir}}
\algrenewcommand\algorithmicuntil{\textbf{hasta que}}

% ---------- Código ----------
\lstset{language=Python,basicstyle=\ttfamily\footnotesize,keywordstyle=\color{azul}\bfseries,
  commentstyle=\color{verde}\itshape,stringstyle=\color{rojo},showstringspaces=false,
  frame=single,framerule=0.3pt,rulecolor=\color{gris!60},breaklines=true,columns=fullflexible,
  keepspaces=true,xleftmargin=4pt,xrightmargin=4pt,aboveskip=8pt,belowskip=8pt,
  literate={á}{{\'a}}1 {é}{{\'e}}1 {í}{{\'i}}1 {ó}{{\'o}}1 {ú}{{\'u}}1 {ñ}{{\~n}}1}

% ---------- Encabezados ----------
\pagestyle{fancy}
\fancyhf{}
\fancyhead[L]{\small\nouppercase{\leftmark}}
\fancyhead[R]{\small\thepage}
\renewcommand{\headrulewidth}{0.3pt}
\fancypagestyle{plain}{\fancyhf{}\fancyfoot[C]{\small\thepage}\renewcommand{\headrulewidth}{0pt}}

% ---------- Macros ----------
\newcommand{\R}{\mathbb{R}}
\newcommand{\E}{\mathbb{E}}
\newcommand{\Prob}{\mathbb{P}}
\newcommand{\cO}{\mathcal{O}}
\newcommand{\cE}{\mathcal{E}}
\newcommand{\cK}{\mathcal{K}}
\newcommand{\cC}{\mathcal{C}}
\newcommand{\cN}{\mathcal{N}}
\newcommand{\cW}{\mathcal{W}}
\newcommand{\T}{^{\top}}
\newcommand{\eps}{\varepsilon}
\newcommand{\norm}[1]{\left\lVert #1\right\rVert}
\newcommand{\nrm}[1]{\lVert #1\rVert}
\newcommand{\abs}[1]{\left\lvert #1\right\rvert}
\newcommand{\ip}[2]{\langle #1,\, #2\rangle}
\DeclareMathOperator{\tr}{tr}
\DeclareMathOperator{\diag}{diag}
\DeclareMathOperator{\rank}{rango}
\DeclareMathOperator{\nnz}{nnz}
\DeclareMathOperator{\dom}{dom}
\DeclareMathOperator{\col}{col}
\DeclareMathOperator*{\argmin}{arg\,min}
\DeclareMathOperator{\cvar}{CVaR}
\newcommand{\lamf}{\lambda_f}
\newcommand{\lamt}{\widetilde{\lambda}_f}
\newcommand{\vne}{v_{\mathrm{ne}}}
\newcommand{\vnsk}{v_{\mathrm{nsk}}}
\newcommand{\xnsk}{x_{\mathrm{nsk}}}
\newcommand{\dmu}{d_{\mu}}
\newcommand{\dmub}{\overline{d}_{\mu}}
\newcommand{\LH}{L_{H}}
\newcommand{\PW}{\textsc{pw}}
\newcommand{\LWP}{\textsc{lwp}}

\setlength{\parskip}{3pt plus 1pt}
\setlist{itemsep=2pt,topsep=4pt}
\allowdisplaybreaks[2]
\emergencystretch=1.5em
\setlength{\belowcaptionskip}{6pt}
\begin{document}
\frontmatter
\begin{titlepage}
\centering
\vspace*{2.2cm}
{\sffamily\Large\color{gris} Documento proyecto\\[2pt] Álgebra lineal numérica y optimización numérica\par}
\vspace{1.6cm}
{\sffamily\bfseries\Huge\color{azul} Newton Sketch adaptativo\par}
\vspace{0.6cm}
{\sffamily\LARGE Matrices de \emph{sketch}, subespacios encajados,\\[4pt] y autoconcordancia\par}
\vspace{1.6cm}
\begin{tikzpicture}[>=Stealth,scale=0.95]
  % A alta
  \draw[fill=azul!18,draw=azul!80] (0,0) rectangle (1.3,4.2);
  \node at (0.65,2.1) {$A$};
  \node[below] at (0.65,0) {\footnotesize $n\times d$};
  % S ancha
  \draw[fill=ocre!25,draw=ocre!90] (-4.4,2.9) rectangle (-0.6,4.2);
  \node at (-2.5,3.55) {$S$};
  \node[below] at (-2.5,2.9) {\footnotesize $m\times n$};
  \node at (-0.3,3.55) {$\cdot$};
  \draw[->,thick] (1.7,2.1) -- (3.3,2.1);
  % SA
  \draw[fill=verde!20,draw=verde!80] (3.7,1.45) rectangle (5.0,2.75);
  \node at (4.35,2.1) {$SA$};
  \node[below] at (4.35,1.45) {\footnotesize $m\times d$};
  \draw[->,thick] (5.4,2.1) -- (7.0,2.1);
  \node[right,align=left] at (7.1,2.1) {$A\T S\T S A \;\approx\; A\T A$\\[4pt]\footnotesize con $m\ll n$};
\end{tikzpicture}
\vfill
{\large Documento con teoría, demostraciones,\\ experimentos numéricos reproducibles \par}
\vspace{1cm}
{\large Autores: Diego Alejandro Arévalo Arias y Yeferson Arlei Getial Cuestas \par}
{\large Noviembre de 2026\par}
\end{titlepage}

\tableofcontents
\mainmatter
\chapter{Introducción y mapa del documento}\label{cap:intro}

\section{¿De qué trata este documento?}

Muchos problemas de optimización a gran escala tienen la forma
\begin{equation}\label{eq:problema-general}
  \min_{x\in\R^d}\; f(x), \qquad\text{con}\qquad \nabla^2 f(x) \;=\; B(x)\T B(x),\quad B(x)\in\R^{n\times d},\quad n\gg d,
\end{equation}
es decir, el hessiano se escribe como el producto de una matriz \emph{alta} por su traspuesta.
Así ocurre en mínimos cuadrados ($B=A$), en regresión logística y demás modelos lineales generalizados
($B=D(x)^{1/2}A$), en programación lineal y cuadrática resueltas con barrera logarítmica
($B=\diag(1/s_i)\,A$, con $s$ el vector de holguras) y, como veremos al final, en el despacho económico de
sistemas eléctricos con muchos escenarios o muchas restricciones de red.

El método de Newton resuelve estos problemas en muy pocas iteraciones, pero cada iteración cuesta
$\cO(nd^2)$ operaciones solo para formar $B\T B$. El descenso de gradiente cuesta $\cO(nd)$ por iteración,
pero su número de iteraciones crece con el número de condición $\kappa=L/\mu$. El \textbf{Newton Sketch}
(Pilanci y Wainwright, 2017) busca lo mejor de ambos mundos: sustituye $B$ por una versión comprimida
$SB\in\R^{m\times d}$, con $S$ una matriz aleatoria de $m\ll n$ filas, y da el paso de Newton con el hessiano
aproximado $(SB)\T(SB)$. El \textbf{Newton Sketch adaptativo} (Lacotte, Wang y Pilanci, 2021) muestra que,
cuando el objetivo tiene una parte fuertemente convexa, el tamaño $m$ puede bajar hasta la \emph{dimensión
efectiva} del hessiano ---que puede ser muy inferior a $d$--- y que dicho tamaño se puede encontrar
automáticamente, duplicándolo solo cuando el progreso observado es insuficiente.

Para entender por qué esto funciona hacen falta tres bloques de teoría que provienen de dos materias distintas:
\begin{enumerate}[label=(\roman*)]
  \item \textbf{Álgebra lineal numérica aleatorizada}: qué matrices $S$ existen, cuánto cuesta aplicarlas y
        por qué $A\T S\T S A\approx A\T A$ (teoría de \emph{subespacios encajados}, en inglés
        \emph{subspace embeddings}).
  \item \textbf{Optimización convexa}: funciones $\mu$-fuertemente convexas y $L$-suaves, método de Newton y,
        sobre todo, \textbf{funciones autoconcordantes}, que permiten analizar Newton sin constantes
        dependientes del problema.
  \item \textbf{Modelado}: cómo aparecen estas estructuras en un problema real, aquí el \textbf{despacho
        económico} (o reparto económico de carga) y sus versiones estocásticas con precios fluctuantes y con
        aportes hídricos gobernados por el fenómeno de El Niño.
\end{enumerate}

\begin{ideaclave}[La idea en una línea]
Un subespacio encajado garantiza $(1-\eps)\,B\T B \preceq B\T S\T S B \preceq (1+\eps)\,B\T B$.
La autoconcordancia garantiza que Newton tolera ese error multiplicativo sin perder su convergencia rápida,
con constantes universales. Juntas dan un método de segundo orden de coste casi lineal en el tamaño de los datos.
\end{ideaclave}

\section{Mapa de dependencias}

\begin{center}
\begin{tikzpicture}[>=Stealth,node distance=9mm and 12mm,
  caja/.style={draw,rounded corners=2pt,align=center,font=\footnotesize,inner sep=4pt,minimum width=3.3cm},
  opt/.style={caja,fill=azul!8,draw=azul!70},
  aln/.style={caja,fill=verde!8,draw=verde!70},
  ns/.style={caja,fill=ocre!12,draw=ocre!90},
  ap/.style={caja,fill=rojo!6,draw=rojo!70}]
  \node[opt] (c2) {Cap.~\ref{cap:opt}: $\mu$, $L$,\\ gradiente y Newton};
  \node[opt,below=of c2] (c3) {Cap.~\ref{cap:sc}: funciones\\ autoconcordantes};
  \node[aln,right=12mm of c2] (c4) {Cap.~\ref{cap:sketch}: matrices\\ de \emph{sketch}};
  \node[aln,below=of c4] (c5) {Cap.~\ref{cap:embed}: subespacios\\ encajados};
  \node[aln,right=7mm of c5] (c6) {Cap.~\ref{cap:ls}--\ref{cap:lin}: mínimos\\ cuadrados y\\ sistemas lineales};
  \node[ns,below=12mm of $(c3)!0.5!(c5)$] (c8) {Cap.~\ref{cap:ns}: Newton Sketch};
  \node[ns,below=of c8] (c9) {Cap.~\ref{cap:ans}: Experimentos};
  \draw[->] (c2) -- (c3);
  \draw[->] (c4) -- (c5);
  \draw[->] (c5) -- (c6);
  \draw[->] (c3) -- (c8);
  \draw[->] (c5) -- (c8);
  \draw[->] (c6.south) |- (c8.east);
  \draw[->] (c8) -- (c9);

\end{tikzpicture}
\end{center}

Los capítulos \ref{cap:opt}--\ref{cap:sc} (azul) son de optimización; los capítulos
\ref{cap:sketch}--\ref{cap:lin} (verde) son de álgebra lineal numérica; ambos confluyen en los capítulos
\ref{cap:ns}--\ref{cap:ans}. Los capítulos \ref{cap:de}--\ref{cap:nino} son la aplicación.



\section{Notación}\label{sec:notacion}

\begin{center}
\small
\begin{tabular}{@{}ll@{}}
\toprule
Símbolo & Significado \\
\midrule
$A\in\R^{n\times d}$ & matriz de datos, alta: $n\gg d$; filas $a_i\T$ \\
$S\in\R^{m\times n}$ & matriz de \emph{sketch}; $m$ es el tamaño del \emph{sketch} \\
$U\in\R^{n\times d}$ & base ortonormal del espacio columna de $A$ \\
$\ell_i$ & puntuación de palanca (\emph{leverage score}) de la fila $i$ \\
$\eps$ & distorsión del encaje, $\eps\in(0,1)$ \\
$p$ & probabilidad de fallo de un evento aleatorio \\
$\delta$ & tolerancia de optimización: se busca $f(x)-f(x^\star)\le\delta$ \\
$\mu,\;L,\;\kappa=L/\mu$ & constantes de convexidad fuerte y de suavidad; número de condición \\
$\LH$ & constante de Lipschitz del hessiano \\
$M$ & constante de autoconcordancia ($M=1$: forma estándar) \\
$H(x)=\nabla^2 f(x)$, $H_S(x)$ & hessiano exacto y hessiano con \emph{sketch} \\
$\nrm{h}_x=(h\T H(x)h)^{1/2}$ & norma local en $x$ \\
$\lamf(x)$, $\lamt(x)$ & decremento de Newton exacto y aproximado \\
$\vne$, $\vnsk$ & dirección de Newton exacta y con \emph{sketch} \\
$\dmu(x)$, $\dmub$ & dimensión efectiva local y global \\
$a,\,b$ & parámetros de la búsqueda lineal con retroceso \\
$\theta$ & parámetro del método de barrera \\
$\nnz(A)$ & número de entradas no nulas de $A$ \\
$X\preceq Y$ & orden de Loewner: $Y-X$ es semidefinida positiva \\
$a\lesssim b$, $a\asymp b$ & desigualdad e igualdad salvo constantes universales \\
\bottomrule
\end{tabular}
\end{center}

\begin{cuidado}[Tres convenciones que difieren de parte de la bibliografía]
\begin{itemize}
  \item En la bibliografía de \emph{sketching} se suele llamar $\delta$ a la probabilidad de fallo. Aquí se
        llama $p$, para reservar $\delta$ a la tolerancia de optimización, como hacen los artículos de Newton Sketch.
\end{itemize}
\end{cuidado}
\part{Optimización: de las constantes del problema a las constantes universales}
\chapter{Convexidad fuerte, suavidad y los métodos clásicos}\label{cap:opt}

Este capítulo fija el lenguaje con el que se mide habitualmente la dificultad de un problema convexo: las
constantes $\mu$ y $L$. Sirve de punto de comparación para todo lo demás: el análisis clásico de Newton
depende de esas constantes; el análisis autoconcordante del capítulo~\ref{cap:sc} no.

\section{Funciones \texorpdfstring{$\mu$}{mu}-fuertemente convexas y \texorpdfstring{$L$}{L}-suaves}

\begin{definicion}[Convexidad fuerte y suavidad]\label{def:mu-L}
Sea $f:\R^d\to\R$ diferenciable.
\begin{enumerate}[label=(\alph*)]
  \item $f$ es \textbf{$\mu$-fuertemente convexa} ($\mu>0$) si, para todo $x,y$,
  \[ f(y)\;\ge\; f(x)+\ip{\nabla f(x)}{y-x}+\frac{\mu}{2}\nrm{y-x}_2^2 . \]
  \item $f$ es \textbf{$L$-suave} si su gradiente es $L$-Lipschitz: $\nrm{\nabla f(x)-\nabla f(y)}_2\le L\nrm{x-y}_2$.
        Si además es convexa, esto equivale a
  \[ f(y)\;\le\; f(x)+\ip{\nabla f(x)}{y-x}+\frac{L}{2}\nrm{y-x}_2^2 . \]
\end{enumerate}
El cociente $\kappa=L/\mu\ge 1$ es el \textbf{número de condición} de $f$.
\end{definicion}

Geométricamente, $f$ queda atrapada entre dos paraboloides: uno de curvatura $\mu$ por debajo y otro de
curvatura $L$ por encima. Para funciones dos veces diferenciables la definición equivale a una cota espectral
uniforme del hessiano:
\begin{equation}\label{eq:mu-L-hess}
  \mu I \;\preceq\; \nabla^2 f(x)\;\preceq\; L I \qquad\text{para todo } x .
\end{equation}

\begin{proposicion}[Consecuencias básicas]\label{prop:mu-L-consec}
Si $f$ es $\mu$-fuertemente convexa y $L$-suave, con minimizador $x^\star$ y valor óptimo $f^\star$, entonces para todo $x$:
\begin{align}
  \frac{\mu}{2}\nrm{x-x^\star}_2^2 \;&\le\; f(x)-f^\star \;\le\; \frac{L}{2}\nrm{x-x^\star}_2^2, \label{eq:cota-dist}\\
  \frac{1}{2L}\nrm{\nabla f(x)}_2^2 \;&\le\; f(x)-f^\star \;\le\; \frac{1}{2\mu}\nrm{\nabla f(x)}_2^2. \label{eq:PL}
\end{align}
\end{proposicion}

\begin{proof}
\eqref{eq:cota-dist} es la definición con $x\leftarrow x^\star$ (donde $\nabla f(x^\star)=0$). Para la cota
superior de \eqref{eq:PL} se minimiza en $y$ el lado derecho de la desigualdad de convexidad fuerte: el mínimo
se alcanza en $y=x-\nabla f(x)/\mu$ y vale $f(x)-\nrm{\nabla f(x)}^2/(2\mu)$; como la desigualdad vale para todo $y$,
en particular $f^\star\ge f(x)-\nrm{\nabla f(x)}^2/(2\mu)$. La cota inferior sale de la suavidad evaluada en
$y=x-\nabla f(x)/L$: $f^\star\le f(y)\le f(x)-\nrm{\nabla f(x)}^2/(2L)$.
\end{proof}

La desigualdad derecha de \eqref{eq:PL} se conoce como condición de Polyak--Łojasiewicz y es lo único que se
necesita para la convergencia lineal del descenso de gradiente.

\begin{ejemplo}[Mínimos cuadrados]\label{ej:ls-mu-L}
Para $f(x)=\tfrac12\nrm{Ax-b}_2^2$ se tiene $\nabla^2 f=A\T A$, de modo que $\mu=\sigma_{\min}(A)^2$,
$L=\sigma_{\max}(A)^2$ y $\kappa(f)=\kappa(A)^2$. Una matriz con $\kappa(A)=10^3$, nada excepcional, da un
problema con $\kappa=10^6$.
\end{ejemplo}

\begin{ejemplo}[Regresión logística regularizada]\label{ej:logistica-mu-L}
Para $f(x)=\sum_{i=1}^n\log(1+e^{-y_i a_i\T x})+\tfrac{\mu}{2}\nrm{x}_2^2$ se tiene
$\nabla^2 f(x)=A\T D(x)A+\mu I$, con $D(x)$ diagonal de entradas $\sigma_i(1-\sigma_i)\in(0,\tfrac14]$. Por tanto
$f$ es $\mu$-fuertemente convexa y $L$-suave con $L=\tfrac14\sigma_{\max}(A)^2+\mu$. Con regularización débil,
$\kappa$ es enorme: en el experimento del capítulo~\ref{cap:ans} vale $4\cdot 10^5$.
\end{ejemplo}

\section{Descenso de gradiente: la factura del número de condición}

\begin{teorema}[Convergencia lineal del descenso de gradiente]\label{teo:gd}
Sea $f$ $\mu$-fuertemente convexa y $L$-suave. La iteración $x_{t+1}=x_t-\tfrac1L\nabla f(x_t)$ satisface
\[ f(x_t)-f^\star\;\le\;\Bigl(1-\frac{\mu}{L}\Bigr)^{t}\,\bigl(f(x_0)-f^\star\bigr). \]
En consecuencia bastan $t\ge\kappa\log\bigl((f(x_0)-f^\star)/\delta\bigr)$ iteraciones para alcanzar precisión $\delta$.
\end{teorema}

\begin{proof}
Por suavidad, $f(x_{t+1})\le f(x_t)-\tfrac{1}{2L}\nrm{\nabla f(x_t)}^2$. Por \eqref{eq:PL},
$\nrm{\nabla f(x_t)}^2\ge 2\mu\,(f(x_t)-f^\star)$. Restando $f^\star$ queda
$f(x_{t+1})-f^\star\le(1-\mu/L)(f(x_t)-f^\star)$.
\end{proof}

El método acelerado de Nesterov mejora la dependencia a $\sqrt{\kappa}\,\log(1/\delta)$ iteraciones, que es
óptima para métodos de primer orden. En ambos casos el coste por iteración es el de un gradiente,
$\cO(nd)$ en problemas del tipo~\eqref{eq:problema-general}, pero el número de iteraciones depende de $\kappa$.

\section{El método de Newton}

El paso de Newton minimiza el modelo cuadrático de $f$ en torno a $x$:
\begin{equation}\label{eq:newton-modelo}
  \vne\;=\;\argmin_{v}\Bigl\{\ip{\nabla f(x)}{v}+\tfrac12\, v\T\nabla^2 f(x)\,v\Bigr\}\;=\;-\nabla^2 f(x)^{-1}\nabla f(x),
  \qquad x^{+}=x+s\,\vne ,
\end{equation}
con $s\in(0,1]$ elegido por \textbf{búsqueda lineal con retroceso} de parámetros $a\in(0,\tfrac12)$, $b\in(0,1)$:
se parte de $s=1$ y se multiplica por $b$ hasta que $f(x+s\vne)\le f(x)+a\,s\,\ip{\nabla f(x)}{\vne}$ (condición de Armijo).

\begin{proposicion}[Invariancia afín]\label{prop:afin}
Sea $T$ invertible y $\tilde f(y)=f(Ty)$. Si $x_0=Ty_0$, las iteraciones de Newton (con la misma búsqueda lineal)
aplicadas a $f$ y a $\tilde f$ satisfacen $x_t=Ty_t$ para todo $t$.
\end{proposicion}

\begin{proof}
$\nabla\tilde f(y)=T\T\nabla f(Ty)$ y $\nabla^2\tilde f(y)=T\T\nabla^2 f(Ty)T$, luego el paso de Newton para
$\tilde f$ es $-T^{-1}\nabla^2 f(Ty)^{-1}\nabla f(Ty)=T^{-1}\vne$. Además $\ip{\nabla \tilde f}{T^{-1}\vne}=\ip{\nabla f}{\vne}$
y $\tilde f(y+sT^{-1}\vne)=f(x+s\vne)$, así que la búsqueda lineal acepta el mismo $s$.
\end{proof}

La invariancia afín es la razón profunda por la que Newton es insensible al mal condicionamiento: un cambio de
variables lineal puede convertir $\kappa=10^6$ en $\kappa=1$ y el método produce exactamente los mismos
iterados. El descenso de gradiente no tiene esta propiedad.

\subsection{El análisis clásico y sus defectos}

\begin{teorema}[Newton con retroceso; análisis clásico {\cite[§9.5.3]{BV}}]\label{teo:newton-clasico}
Sea $f$ $\mu$-fuertemente convexa, $L$-suave y con hessiano $\LH$-Lipschitz:
$\nrm{\nabla^2 f(x)-\nabla^2 f(y)}_2\le\LH\nrm{x-y}_2$. Existen
\[ \eta=\min\{1,\,3(1-2a)\}\,\frac{\mu^2}{\LH}, \qquad \gamma=a\,b\,\eta^2\,\frac{\mu}{L^2}, \]
tales que:
\begin{enumerate}[label=(\roman*)]
  \item \textbf{Fase amortiguada.} Si $\nrm{\nabla f(x_t)}_2\ge\eta$, entonces $f(x_{t+1})-f(x_t)\le-\gamma$.
  \item \textbf{Fase cuadrática.} Si $\nrm{\nabla f(x_t)}_2<\eta$, la búsqueda lineal acepta $s=1$ y
  \[ \frac{\LH}{2\mu^2}\nrm{\nabla f(x_{t+1})}_2\;\le\;\Bigl(\frac{\LH}{2\mu^2}\nrm{\nabla f(x_t)}_2\Bigr)^{2}. \]
\end{enumerate}
El número total de iteraciones para alcanzar $f(x)-f^\star\le\delta$ es a lo sumo
\[ \frac{f(x_0)-f^\star}{\gamma}\;+\;\log_2\log_2\frac{\eps_0}{\delta},\qquad \eps_0=\frac{2\mu^3}{\LH^2}. \]
\end{teorema}

El término $\log_2\log_2(\eps_0/\delta)$ es, a efectos prácticos, una constante: vale $6$ para una precisión
relativa de $10^{-19}$. Este es el sentido de \emph{convergencia cuadrática}: el número de cifras correctas se
duplica en cada iteración.

\begin{cuidado}[Lo que este teorema tiene de insatisfactorio]
\begin{enumerate}
  \item La cota depende de $\mu$, $L$ y $\LH$, que casi nunca se conocen, y a través de $\gamma\propto\mu^5/(L^2\LH^2)$
        es pésima para problemas mal condicionados; en la práctica Newton se comporta muchísimo mejor.
  \item La cota cambia si se hace un cambio de variables lineal, pero el algoritmo no (proposición~\ref{prop:afin}).
        Un análisis en norma euclídea no puede capturar un método que es invariante afín.
\end{enumerate}
La teoría de funciones autoconcordantes (capítulo~\ref{cap:sc}) resuelve ambos problemas, y es la que usan
los teoremas de Newton Sketch.
\end{cuidado}

\subsection{El coste de una iteración y la raíz cuadrada del hessiano}\label{sec:raiz-hess}

En los problemas que nos interesan el hessiano tiene una \emph{raíz cuadrada} explícita:
una matriz $B(x)\in\R^{n\times d}$ con $B(x)\T B(x)=\nabla^2 f(x)$. (No es la raíz simétrica $d\times d$: es
rectangular y se obtiene gratis de la estructura del problema.)

\begin{center}
\small
\begin{tabular}{@{}lll@{}}
\toprule
Problema & $f(x)$ & $B(x)=\nabla^2 f(x)^{1/2}$ \\
\midrule
Mínimos cuadrados & $\tfrac12\nrm{Ax-b}_2^2$ & $A$ \\
Modelo lineal generalizado & $\sum_i\varphi_i(a_i\T x)$ & $\diag\bigl(\varphi_i''(a_i\T x)^{1/2}\bigr)A$ \\
Regresión logística & $\sum_i\log(1+e^{-y_ia_i\T x})$ & $\diag\bigl((\sigma_i(1-\sigma_i))^{1/2}\bigr)A$ \\
Barrera logarítmica & $-\sum_i\log(b_i-a_i\T x)$ & $\diag\bigl(1/(b_i-a_i\T x)\bigr)A$ \\
\bottomrule
\end{tabular}
\end{center}

Con esta estructura, una iteración de Newton cuesta $\cO(nd)$ para formar $B$ y el gradiente,
$\cO(nd^2)$ para formar $B\T B$ y $\cO(d^3)$ para resolver el sistema. Cuando $n\gg d$, \textbf{el cuello de
botella es el producto $B\T B$}. Ahí es exactamente donde actúa el \emph{sketching}.

\begin{resumen}
\begin{itemize}
  \item $\mu$ y $L$ acotan la curvatura por abajo y por arriba; $\kappa=L/\mu$ mide la dificultad para los métodos de primer orden.
  \item Descenso de gradiente: $\cO(nd)$ por iteración, $\kappa\log(1/\delta)$ iteraciones.
  \item Newton: $\cO(nd^2)$ por iteración, número de iteraciones prácticamente independiente de $\kappa$; es invariante afín.
  \item El análisis clásico de Newton depende de $\mu$, $L$, $\LH$ y no refleja su comportamiento real.
  \item El hessiano de los problemas de interés es $B\T B$ con $B$ alta: formar ese producto es lo caro.
\end{itemize}
\end{resumen}


\chapter{Funciones autoconcordantes}\label{cap:sc}

La autoconcordancia, introducida por Nesterov y Nemirovski, es la condición de regularidad \emph{adaptada al
método de Newton}. En lugar de acotar el hessiano por constantes fijas ($\mu$ y $L$), acota \emph{cuánto
puede cambiar el hessiano} medido en la métrica que el propio hessiano define. Es la hipótesis central de
los teoremas de Newton Sketch.

\section{Definición y primeros ejemplos}

\begin{definicion}[Autoconcordancia]\label{def:sc}
Una función convexa $\varphi:\R\to\R\cup\{+\infty\}$, tres veces derivable en su dominio abierto, es
\textbf{autoconcordante con constante $M\ge0$} si
\begin{equation}\label{eq:sc-1d}
  \abs{\varphi'''(t)}\;\le\;2M\,\varphi''(t)^{3/2}\qquad\text{para todo } t\in\dom\varphi .
\end{equation}
Una función convexa $f:\R^d\to\R\cup\{+\infty\}$ es autoconcordante con constante $M$ si lo es su restricción
a cada recta, $t\mapsto f(x+th)$. Equivalentemente,
\begin{equation}\label{eq:sc-nd}
  \abs{D^3f(x)[h,h,h]}\;\le\;2M\,\bigl(h\T\nabla^2 f(x)\,h\bigr)^{3/2}\qquad\text{para todo } x\in\dom f,\ h\in\R^d .
\end{equation}
Si $M=1$ se dice \textbf{autoconcordante estándar}, o simplemente autoconcordante.
\end{definicion}

El exponente $3/2$ no es arbitrario: es el único que hace la condición invariante bajo cambios de variable
afines $x\mapsto Tx+c$ (ambos lados de \eqref{eq:sc-nd} se transforman igual). El factor $2$ es una
normalización que hace que $-\log t$ tenga constante exactamente $1$.

\begin{ejemplo}[Funciones autoconcordantes]\label{ej:sc}
\begin{enumerate}[label=(\alph*)]
  \item \textbf{Lineales y cuadráticas convexas}: $\varphi'''=0$; son autoconcordantes con $M=0$.
  \item \textbf{Logaritmo negativo}: $\varphi(t)=-\log t$ en $t>0$. $\varphi''=1/t^2$, $\varphi'''=-2/t^3$, de modo
        que $\abs{\varphi'''}=2(\varphi'')^{3/2}$: se cumple \eqref{eq:sc-1d} con igualdad y $M=1$.
  \item \textbf{Barrera logarítmica de un poliedro}: $f(x)=-\sum_{i=1}^n\log(b_i-a_i\T x)$; es suma de
        composiciones afines de (b), luego autoconcordante estándar (proposición~\ref{prop:sc-calculo}).
  \item \textbf{Entropía más barrera}: $\varphi(t)=t\log t-\log t$ en $t>0$.
  \item \textbf{Barrera del cono semidefinido}: $f(X)=-\log\det X$ sobre $X\succ0$.
\end{enumerate}
\end{ejemplo}

\begin{ejemplo}[Funciones que \emph{no} lo son]\label{ej:no-sc}
\begin{enumerate}[label=(\alph*)]
  \item $\varphi(t)=e^{t}$: $\varphi'''/(\varphi'')^{3/2}=e^{-t/2}\to\infty$ cuando $t\to-\infty$.
  \item La pérdida logística $\varphi(t)=\log(1+e^{-t})$: con $\sigma=1/(1+e^{t})$ se tiene
        $\varphi''=\sigma(1-\sigma)$ y $\abs{\varphi'''}=\sigma(1-\sigma)\abs{1-2\sigma}$, de modo que
        $\abs{\varphi'''}/(\varphi'')^{3/2}=\abs{1-2\sigma}/\sqrt{\sigma(1-\sigma)}\to\infty$ cuando $\abs{t}\to\infty$.
        Cumple, en cambio, $\abs{\varphi'''}\le\varphi''$ (autoconcordancia \emph{generalizada}), y esto basta
        para recuperar la autoconcordancia al añadir una regularización (proposición~\ref{prop:logistica-sc}).
\end{enumerate}
\end{ejemplo}

\begin{proposicion}[Cálculo de funciones autoconcordantes]\label{prop:sc-calculo}
\begin{enumerate}[label=(\alph*)]
  \item \textbf{Composición afín}: si $f$ es autoconcordante con constante $M$, también lo es $x\mapsto f(Tx+c)$.
  \item \textbf{Suma}: si $f_1$ y $f_2$ lo son con constantes $M_1$, $M_2$, entonces $f_1+f_2$ lo es con constante $\max\{M_1,M_2\}$.
  \item \textbf{Escalado}: si $f$ lo es con constante $M$ y $c>0$, entonces $cf$ lo es con constante $M/\sqrt{c}$.
        En particular, $M^2 f$ es autoconcordante estándar.
\end{enumerate}
\end{proposicion}

\begin{proof}
(a) Basta restringir a rectas. (b) En una recta, con $M=\max\{M_1,M_2\}$,
$\abs{\varphi_1'''+\varphi_2'''}\le 2M\bigl((\varphi_1'')^{3/2}+(\varphi_2'')^{3/2}\bigr)\le 2M(\varphi_1''+\varphi_2'')^{3/2}$,
porque $u^{3/2}+w^{3/2}\le(u+w)^{3/2}$ para $u,w\ge0$. (c) $\abs{c\varphi'''}\le 2Mc(\varphi'')^{3/2}=2Mc^{-1/2}(c\varphi'')^{3/2}$.
\end{proof}

\begin{cuidado}[El escalado importa]
El apartado (c) dice que multiplicar por una constante \emph{pequeña} empeora la constante de
autoconcordancia. El paso de Newton no cambia al escalar $f$, pero los umbrales de los teoremas (enunciados
para $M=1$) sí se refieren a la escala estándar: el decremento de Newton de $M^2f$ es $M$ veces el de $f$.
\end{cuidado}

\section{La geometría local: norma local y elipsoide de Dikin}

Sea $f$ autoconcordante estándar con $\nabla^2 f(x)\succ0$. El hessiano define en cada punto una norma y su dual:
\[ \nrm{h}_x=\bigl(h\T\nabla^2 f(x)\,h\bigr)^{1/2},\qquad \nrm{g}_x^{*}=\bigl(g\T\nabla^2 f(x)^{-1}g\bigr)^{1/2}. \]
El \textbf{elipsoide de Dikin} de radio $r$ es $W(x,r)=\{y:\nrm{y-x}_x<r\}$.

El hecho básico del que se deduce todo lo demás es un cálculo de una línea. Si $\varphi$ es autoconcordante
estándar en una variable y $\varphi''>0$,
\begin{equation}\label{eq:sc-clave}
  \abs{\frac{d}{dt}\Bigl(\varphi''(t)^{-1/2}\Bigr)}\;=\;\abs{\frac{\varphi'''(t)}{2\,\varphi''(t)^{3/2}}}\;\le\;1 .
\end{equation}
Es decir: \emph{la función $\varphi''^{-1/2}$ es $1$-Lipschitz}. Integrando entre $0$ y $t$ se obtiene
$\varphi''(0)^{-1/2}-t\le\varphi''(t)^{-1/2}\le\varphi''(0)^{-1/2}+t$, o sea, con $c_0=\varphi''(0)^{1/2}$,
\begin{equation}\label{eq:sc-hess-1d}
  \frac{\varphi''(0)}{(1+c_0t)^2}\;\le\;\varphi''(t)\;\le\;\frac{\varphi''(0)}{(1-c_0t)^2},
  \qquad 0\le t<1/c_0 .
\end{equation}
Aplicando esto a $\varphi(t)=f(x+th)$ e integrando una y dos veces se obtienen las cotas siguientes.

\begin{teorema}[Cotas fundamentales {\cite[§5.1]{Nesterov}}]\label{teo:sc-cotas}
Sea $f$ autoconcordante estándar, $x\in\dom f$ y $r=\nrm{y-x}_x$. Con
$\omega(\tau)=\tau-\log(1+\tau)$ y $\omega_*(\tau)=-\tau-\log(1-\tau)$:
\begin{enumerate}[label=(\roman*)]
  \item \textbf{Elipsoide de Dikin}: si $r<1$ entonces $y\in\dom f$; es decir, $W(x,1)\subseteq\dom f$.
  \item \textbf{Hessiano}: si $r<1$,
  \begin{equation}\label{eq:sc-hess}
    (1-r)^2\,\nabla^2 f(x)\;\preceq\;\nabla^2 f(y)\;\preceq\;\frac{1}{(1-r)^2}\,\nabla^2 f(x).
  \end{equation}
  \item \textbf{Cota inferior}: para todo $y\in\dom f$, $\;f(y)\ge f(x)+\ip{\nabla f(x)}{y-x}+\omega(r)$.
  \item \textbf{Cota superior}: si $r<1$, $\;f(y)\le f(x)+\ip{\nabla f(x)}{y-x}+\omega_*(r)$.
\end{enumerate}
\end{teorema}

Como $\omega(\tau)\approx\tau^2/2\approx\omega_*(\tau)$ para $\tau$ pequeño, (iii) y (iv) dicen que cerca de $x$ la
función es, en la norma local, casi exactamente una cuadrática de curvatura~$1$. Lejos de $x$ la cota
inferior crece solo linealmente ($\omega(\tau)\sim\tau$) y la superior explota en $\tau=1$: nada se garantiza
fuera del elipsoide de Dikin.

\begin{ideaclave}
Las desigualdades \eqref{eq:sc-hess} son el análogo de $\mu I\preceq\nabla^2 f\preceq LI$, pero con dos
diferencias decisivas: comparan el hessiano \emph{consigo mismo en un punto cercano} (no con la identidad) y
las «constantes» $(1-r)^{\pm2}$ son \emph{universales}: no dependen de la función ni de los datos.
\end{ideaclave}

\section{Decremento de Newton y análisis del método}

\begin{definicion}[Decremento de Newton]\label{def:decremento}
\[ \lamf(x)\;=\;\nrm{\nabla f(x)}_x^{*}\;=\;\bigl(\nabla f(x)\T\nabla^2 f(x)^{-1}\nabla f(x)\bigr)^{1/2}. \]
\end{definicion}

Tres lecturas equivalentes, todas útiles: (1) $\lamf(x)=\nrm{\vne}_x$ es la longitud del paso de Newton en la
norma local; (2) $\lamf(x)^2=-\ip{\nabla f(x)}{\vne}$ es la derivada direccional que usa la condición de Armijo;
(3) $\tfrac12\lamf(x)^2$ es lo que el modelo cuadrático \eqref{eq:newton-modelo} predice que bajará $f$.
Además $\lamf$ es invariante afín, y por eso es el criterio de parada natural.

\begin{proposicion}[El decremento acota el error]\label{prop:gap-lambda}
Si $f$ es autoconcordante estándar, acotada inferiormente, y $\lamf(x)<1$, entonces
\[ f(x)-f^\star\;\le\;\omega_*(\lamf(x))=-\lamf(x)-\log(1-\lamf(x)), \]
y en particular $f(x)-f^\star\le\lamf(x)^2$ siempre que $\lamf(x)\le 0.68$.
\end{proposicion}

\begin{proof}
Por el teorema~\ref{teo:sc-cotas}(iii) y Cauchy--Schwarz en las normas local y dual,
$f(y)\ge f(x)-\lamf(x)\,r+\omega(r)$ con $r=\nrm{y-x}_x$. Minimizando en $r\ge0$ (el mínimo está en
$r=\lamf/(1-\lamf)$) resulta $f(y)\ge f(x)-\omega_*(\lamf(x))$ para todo $y$. La segunda afirmación es la
desigualdad numérica $\omega_*(\tau)\le\tau^2$ en $[0,\,0.68]$.
\end{proof}

\begin{teorema}[Newton amortiguado: descenso garantizado]\label{teo:sc-damped}
Sea $f$ autoconcordante estándar. El paso $x^+=x+\frac{1}{1+\lamf(x)}\,\vne$ pertenece a $\dom f$ y
\[ f(x^+)\;\le\;f(x)-\omega(\lamf(x))=f(x)-\lamf(x)+\log(1+\lamf(x)). \]
\end{teorema}

\begin{proof}
Sea $\lambda=\lamf(x)$. Para $s\lambda<1$, el teorema~\ref{teo:sc-cotas}(iv) con $y=x+s\vne$ da
$f(x+s\vne)\le f(x)-s\lambda^2+\omega_*(s\lambda)$. El lado derecho se minimiza en $s=1/(1+\lambda)$, donde
$s\lambda=\lambda/(1+\lambda)<1$, y vale
$f(x)-\tfrac{\lambda^2}{1+\lambda}-\tfrac{\lambda}{1+\lambda}-\log\bigl(1-\tfrac{\lambda}{1+\lambda}\bigr)=f(x)-\lambda+\log(1+\lambda)$.
\end{proof}

Mientras $\lamf(x)\ge\eta$, cada iteración baja $f$ en al menos la constante $\omega(\eta)$: esta es la
\emph{fase amortiguada}. Cuando $\lamf$ es pequeño empieza la \emph{fase cuadrática}:

\begin{teorema}[Convergencia cuadrática {\cite[§5.2]{Nesterov}}]\label{teo:sc-quad}
Sea $f$ autoconcordante estándar y $\lamf(x)<1$. El paso completo $x^+=x+\vne$ satisface
\[ \lamf(x^+)\;\le\;\Bigl(\frac{\lamf(x)}{1-\lamf(x)}\Bigr)^{2}. \]
En particular, si $\lamf(x)\le\tfrac14$ entonces $\lamf(x^+)\le2\,\lamf(x)^2$.
\end{teorema}

\begin{proof}[Esquema]
Como $\nabla f(x)+\nabla^2 f(x)\vne=0$, se tiene
$\nabla f(x^+)=\int_0^1\bigl(\nabla^2 f(x+\tau\vne)-\nabla^2 f(x)\bigr)\vne\,d\tau$. Por \eqref{eq:sc-hess} con
$r=\tau\lamf(x)$, la diferencia de hessianos está acotada, en la métrica de $\nabla^2 f(x)$, por
$(1-\tau\lambda)^{-2}-1$, con $\lambda=\lamf(x)$; integrando resulta $\nrm{\nabla f(x^+)}_x^{*}\le\lambda^2/(1-\lambda)$. Pasar de la norma dual en
$x$ a la norma dual en $x^+$ cuesta otro factor $(1-\lambda)^{-1}$, de nuevo por \eqref{eq:sc-hess}.
\end{proof}

Combinando ambas fases con la búsqueda lineal con retroceso se obtiene el teorema de complejidad que servirá
de modelo para Newton Sketch.

\begin{teorema}[Complejidad de Newton para funciones autoconcordantes {\cite[§9.6.4]{BV}}]\label{teo:sc-complejidad}
Sea $f$ autoconcordante estándar, estrictamente convexa y acotada inferiormente. El método de Newton con
búsqueda lineal con retroceso de parámetros $(a,b)$ alcanza $f(x)-f^\star\le\delta$ en a lo sumo
\[ \frac{f(x_0)-f^\star}{\gamma}\;+\;\log_2\log_2\frac1\delta \qquad\text{iteraciones},\qquad
   \gamma=\frac{a\,b\,(1-2a)^2}{20-8a}. \]
La transición entre fases ocurre en $\lamf(x)\le\eta=(1-2a)/4$.
\end{teorema}

Compárese con el teorema~\ref{teo:newton-clasico}: \textbf{no aparece ninguna constante del problema}. Solo
intervienen los parámetros de la búsqueda lineal y la distancia inicial al óptimo \emph{en valor de la función}.
Con $a=0.1$, $b=0.5$ se obtiene $\gamma\approx1/600$; la cota es muy pesimista en la práctica (lo habitual son
unas pocas decenas de iteraciones), pero su estructura es la correcta.

\section{Relación con las funciones \texorpdfstring{$\mu$}{mu}-fuertemente convexas y \texorpdfstring{$L$}{L}-suaves}\label{sec:sc-vs-muL}

Son tres clases distintas que se solapan. Conviene tener claro qué implica qué.

\begin{proposicion}[Fuertemente convexa con hessiano Lipschitz $\Rightarrow$ autoconcordante]\label{prop:muL-sc}
Si $f$ es $\mu$-fuertemente convexa y su hessiano es $\LH$-Lipschitz, entonces $f$ es autoconcordante con constante
\[ M=\frac{\LH}{2\,\mu^{3/2}}, \]
y por tanto $\dfrac{\LH^2}{4\mu^3}\,f$ es autoconcordante estándar.
\end{proposicion}

\begin{proof}
Que el hessiano sea $\LH$-Lipschitz implica $\abs{D^3f(x)[h,h,h]}\le\LH\nrm{h}_2^3$. Por convexidad fuerte,
$\nrm{h}_2^2\le\tfrac1\mu\,h\T\nabla^2 f(x)h$. Luego
$\abs{D^3f(x)[h,h,h]}\le\LH\,\mu^{-3/2}(h\T\nabla^2 f(x)h)^{3/2}$, que es \eqref{eq:sc-nd} con $2M=\LH\mu^{-3/2}$.
\end{proof}

\begin{proposicion}[Recíproco parcial]\label{prop:sc-muL}
Si $f$ es autoconcordante con constante $M$ y además $\nabla^2 f(x)\preceq LI$ en un conjunto convexo, entonces
el hessiano es Lipschitz en ese conjunto con constante $\LH=2ML^{3/2}$.
\end{proposicion}

\begin{proof}
$\abs{D^3f(x)[h,h,h]}\le2M(h\T\nabla^2 f(x)h)^{3/2}\le2ML^{3/2}\nrm{h}_2^3$. Para una forma trilineal simétrica el
máximo de $\abs{T[u,v,w]}$ sobre vectores unitarios se alcanza con $u=v=w$, luego la norma de operador de
$D^3f(x)$ es a lo sumo $2ML^{3/2}$, y esa norma acota la constante de Lipschitz de $\nabla^2 f$.
\end{proof}

Así, \emph{dentro} de la clase de funciones con $\mu I\preceq\nabla^2 f\preceq LI$, ser autoconcordante (con
alguna constante) equivale a tener hessiano Lipschitz. Fuera de esa clase ninguna noción implica la otra:

\begin{proposicion}[Ninguna implicación general]\label{prop:contraejemplos}
\begin{enumerate}[label=(\alph*)]
  \item \textbf{Autoconcordante $\not\Rightarrow$ fuertemente convexa ni suave.} $f(t)=-\log t$ en $(0,\infty)$:
        $f''(t)=1/t^2$ tiende a $0$ cuando $t\to\infty$ (no hay $\mu>0$) y a $\infty$ cuando $t\to0$ (no hay $L<\infty$).
        Una función lineal es autoconcordante y ni siquiera es estrictamente convexa.
  \item \textbf{Fuertemente convexa y suave $\not\Rightarrow$ autoconcordante.} $f(t)=\tfrac12t^2+\tfrac{4}{15}\abs{t}^{5/2}$:
        $f''(t)=1+\abs{t}^{1/2}\in[1,2]$ en $[-1,1]$, pero $f'''(t)=\pm\tfrac12\abs{t}^{-1/2}$ no está acotada
        cerca de $0$, así que ningún $M$ finito sirve.
\end{enumerate}
\end{proposicion}

\subsection{La relación conceptual: buen condicionamiento local universal}

La manera más útil de entender la relación es reescribir \eqref{eq:sc-hess} como una afirmación sobre $\mu$ y
$L$ \emph{en otra métrica}.

\begin{proposicion}[Toda función autoconcordante es localmente bien condicionada en su norma local]\label{prop:cond-local}
Sea $f$ autoconcordante estándar, $x\in\dom f$ con $\nabla^2 f(x)\succ0$ y $r\in(0,1)$. En el elipsoide de
Dikin $W(x,r)$, la función $f$ es $(1-r)^2$-fuertemente convexa y $(1-r)^{-2}$-suave \emph{respecto de la
norma $\nrm{\cdot}_x$}. Su número de condición local es, por tanto, a lo sumo $(1-r)^{-4}$: por ejemplo, $16$
en el elipsoide de radio $\tfrac12$, sea cual sea la función.
\end{proposicion}

\begin{proof}
Es \eqref{eq:sc-hess}: en las variables $z=\nabla^2 f(x)^{1/2}y$, la norma $\nrm{\cdot}_x$ es la euclídea y el
hessiano de $f$ queda entre $(1-r)^2I$ y $(1-r)^{-2}I$ en todo el elipsoide.
\end{proof}

\begin{ideaclave}[Tres maneras de mirar lo mismo]
\begin{itemize}
  \item Una función \textbf{$\mu$-fuertemente convexa y $L$-suave} tiene curvatura controlada \emph{globalmente
        y en la métrica euclídea}; el control puede ser malo ($\kappa$ grande).
  \item Una función \textbf{autoconcordante} tiene curvatura controlada \emph{localmente y en su propia
        métrica}; el control es siempre bueno (constantes universales), pero solo vale dentro del elipsoide de Dikin.
  \item El \textbf{método de Newton} es descenso de gradiente en la métrica local: $\vne=-\nabla^2 f(x)^{-1}\nabla f(x)$
        es el gradiente respecto del producto interno $\ip{u}{v}_x=u\T\nabla^2 f(x)v$. Por eso hereda el buen
        condicionamiento local y su complejidad no depende de $\kappa$.
\end{itemize}
\end{ideaclave}

\begin{observacion}[En conjuntos de nivel acotados]\label{obs:nivel}
Si $f$ es autoconcordante y estrictamente convexa y el conjunto de nivel $\{x:f(x)\le f(x_0)\}$ es compacto,
el hessiano alcanza allí un mínimo $\mu>0$ y un máximo $L<\infty$ de sus valores propios extremos: restringida a
ese conjunto, $f$ \emph{es} $\mu$-fuertemente convexa y $L$-suave. Pero $\kappa=L/\mu$ puede ser astronómico
(piénsese en una barrera cerca de la frontera). El descenso de gradiente paga ese $\kappa$; Newton no.
\end{observacion}

\subsection{El caso de la regresión logística regularizada}

La pérdida logística no es autoconcordante (ejemplo~\ref{ej:no-sc}), pero al sumarle un término fuertemente
convexo sí lo es, con una constante que depende de $\mu$. Este es el mecanismo por el que la regresión
logística regularizada entra en el marco de Newton Sketch.

\begin{proposicion}\label{prop:logistica-sc}
Sea $f(x)=\sum_{i=1}^n\varphi_i(a_i\T x)+\tfrac\mu2\nrm{x}_2^2$, con $\varphi_i$ convexas que cumplen
$\abs{\varphi_i'''}\le\varphi_i''$, y sea $R=\max_i\nrm{a_i}_2$. Entonces $f$ es autoconcordante con constante
\[ M=\frac{R}{2\sqrt{\mu}},\qquad\text{y por tanto}\quad \frac{R^2}{4\mu}\,f\ \text{ es autoconcordante estándar.} \]
\end{proposicion}

\begin{proof}
Sea $H=\nabla^2 f(x)=\sum_i\varphi_i''(a_i\T x)a_ia_i\T+\mu I$. Entonces
\[ \abs{D^3f(x)[h,h,h]}=\Bigl|\sum_i\varphi_i'''(a_i\T x)(a_i\T h)^3\Bigr|
   \le\max_i\abs{a_i\T h}\,\sum_i\varphi_i''(a_i\T x)(a_i\T h)^2\le R\nrm{h}_2\;h\T Hh . \]
Como $H\succeq\mu I$, $\nrm{h}_2\le\mu^{-1/2}(h\T Hh)^{1/2}$, y queda
$\abs{D^3f(x)[h,h,h]}\le\tfrac{R}{\sqrt\mu}(h\T Hh)^{3/2}$.
\end{proof}

\section{Barreras y método de punto interior}\label{sec:barrera}

La aplicación más importante de la autoconcordancia es convertir un problema con restricciones en una
sucesión de problemas sin restricciones. Para
\begin{equation}\label{eq:prob-restr}
  \min_x\; c(x)\quad\text{s.a.}\quad a_i\T x\le b_i,\ \ i=1,\dots,n,
\end{equation}
con $c$ convexa (lineal o cuadrática en nuestros ejemplos), se define la \textbf{barrera logarítmica}
$\phi(x)=-\sum_{i=1}^n\log(b_i-a_i\T x)$ y, para cada $\theta>0$, el problema de \emph{centrado}
\begin{equation}\label{eq:centrado}
  x^\star(\theta)=\argmin_x\;F_\theta(x),\qquad F_\theta(x)=\theta\,c(x)+\phi(x).
\end{equation}
La curva $\theta\mapsto x^\star(\theta)$ es la \textbf{trayectoria central}. Se cumple
$c(x^\star(\theta))-c^\star\le n/\theta$: el punto central es $n/\theta$-subóptimo.

\begin{algorithm}[ht]
\caption{Método de barrera (seguimiento de la trayectoria central)}\label{alg:barrera}
\begin{algorithmic}[1]
\Require $x$ estrictamente factible, $\theta_0>0$, factor $\mu_b>1$, tolerancia $\delta$
\State $\theta\gets\theta_0$
\Repeat
  \State \textbf{Centrado:} partiendo de $x$, minimizar $F_\theta$ con Newton (o Newton Sketch); $x\gets x^\star(\theta)$
  \State \textbf{si} $n/\theta\le\delta$ \textbf{entonces devolver} $x$
  \State $\theta\gets\mu_b\,\theta$
\Until{convergencia}
\end{algorithmic}
\end{algorithm}

Si $c$ es lineal o cuadrática convexa, $F_\theta$ es autoconcordante estándar para todo $\theta$
(proposición~\ref{prop:sc-calculo}), de modo que cada centrado se analiza con el
teorema~\ref{teo:sc-complejidad}. El número de centrados es $\lceil\log(n/(\theta_0\delta))/\log\mu_b\rceil$, y
se puede probar que cada uno requiere un número de pasos de Newton acotado en función solo de $n$ y $\mu_b$;
con la elección $\mu_b=1+1/\sqrt n$ el total es $\cO(\sqrt n\log(n/(\theta_0\delta)))$.

El hessiano de la barrera es $\nabla^2\phi(x)=A\T\diag(1/s_i^2)A$ con holguras $s_i=b_i-a_i\T x$, y su raíz
cuadrada es $\diag(1/s_i)A\in\R^{n\times d}$: \textbf{cuando hay muchas más restricciones que variables,
estamos exactamente en la situación \eqref{eq:problema-general}}.

\begin{resumen}
\begin{itemize}
  \item Autoconcordancia: $\abs{\varphi'''}\le2M(\varphi'')^{3/2}$ en cada recta; es invariante afín.
  \item Consecuencia central: $(1-r)^2H(x)\preceq H(y)\preceq(1-r)^{-2}H(x)$ si $r=\nrm{y-x}_x<1$.
  \item El decremento $\lamf$ mide el error ($f-f^\star\le\lamf^2$) y gobierna las dos fases de Newton.
  \item Complejidad de Newton: $(f(x_0)-f^\star)/\gamma+\log_2\log_2(1/\delta)$, sin constantes del problema.
  \item $\mu$-fuerte $+$ hessiano Lipschitz $\Rightarrow$ autoconcordante con $M=\LH/(2\mu^{3/2})$; en general ninguna clase contiene a la otra.
  \item Las barreras logarítmicas son autoconcordantes y su hessiano es $B\T B$ con $B=\diag(1/s)A$.
\end{itemize}
\end{resumen}

\part{Álgebra lineal numérica aleatorizada}
\chapter{Matrices de \emph{sketch}}\label{cap:sketch}

Un \emph{sketch} (boceto, esbozo) de una matriz alta $A\in\R^{n\times d}$ es una matriz mucho más pequeña
$SA\in\R^{m\times d}$, con $m\ll n$, que conserva la información geométrica que nos importa. Este capítulo
describe las matrices $S$ principales: cómo se construyen, cuánto cuesta aplicarlas y por qué funcionan. La
teoría que cuantifica «funcionar» se desarrolla en el capítulo~\ref{cap:embed}.

\section{Qué se le pide a una matriz de \emph{sketch}}

\begin{definicion}[Matriz de \emph{sketch}]\label{def:sketch}
Una matriz de \emph{sketch} es una matriz aleatoria $S\in\R^{m\times n}$, $m\ll n$, normalizada de modo que
\begin{equation}\label{eq:isotropia}
  \E\bigl[S\T S\bigr]=I_n .
\end{equation}
\end{definicion}

La normalización \eqref{eq:isotropia} dice que $S$ conserva las normas \emph{en media}:
$\E\nrm{Sx}_2^2=x\T\E[S\T S]x=\nrm{x}_2^2$ para todo $x$ fijo. Lo que distingue un buen \emph{sketch} de uno malo
es la \emph{concentración}: que $\nrm{Sx}_2^2\approx\nrm{x}_2^2$ con alta probabilidad y, sobre todo, que
ocurra \emph{simultáneamente para todos los $x$ de un subespacio}. Los tres criterios de diseño son:
\begin{enumerate}
  \item \textbf{Tamaño}: cuántas filas $m$ hacen falta para una distorsión $\eps$ dada.
  \item \textbf{Coste de aplicación}: cuántas operaciones requiere calcular $SA$.
  \item \textbf{Olvido de los datos} (\emph{obliviousness}): si $S$ se puede generar sin mirar $A$.
\end{enumerate}

\section{\emph{Sketches} densos: gaussiano y subgaussiano}

\begin{definicion}[\emph{Sketch} gaussiano]
$S=\frac{1}{\sqrt m}G$, donde $G\in\R^{m\times n}$ tiene entradas independientes $N(0,1)$.
\end{definicion}

\paragraph{Cómo funciona.} Cada fila de $SA$ es una combinación lineal aleatoria de \emph{todas} las filas de
$A$. La propiedad que hace del caso gaussiano el más limpio de analizar es la \textbf{invariancia por
rotaciones}: si $U\in\R^{n\times d}$ tiene columnas ortonormales, $GU$ es de nuevo una matriz gaussiana
estándar, de tamaño $m\times d$. Por tanto $SU$ se comporta como una matriz gaussiana \emph{que ya no
depende de $n$}, y sus valores singulares satisfacen (Davidson--Szarek), con probabilidad al menos $1-2e^{-m\tau^2/2}$,
\begin{equation}\label{eq:gauss-sv}
  1-\sqrt{d/m}-\tau\;\le\;\sigma_{\min}(SU)\;\le\;\sigma_{\max}(SU)\;\le\;1+\sqrt{d/m}+\tau .
\end{equation}
La distorsión es, pues, $\eps\approx2\sqrt{d/m}$: para $\eps$ fijo basta $m\asymp d/\eps^2$, sin logaritmos.

\paragraph{Coste.} Formar $SA$ es un producto denso: $\cO(mnd)$. Con $m\asymp d$ esto es $\cO(nd^2)$, lo mismo
que formar $A\T A$. El \emph{sketch} gaussiano es la referencia teórica, no la opción práctica.

\paragraph{Variantes.} Los mismos resultados valen, con otras constantes, para entradas independientes
\emph{subgaussianas} de media cero y varianza $1/m$; por ejemplo signos de Rademacher $\pm1/\sqrt m$.

\section{Transformada de Hadamard aleatorizada submuestreada (SRHT)}

\begin{definicion}[SRHT]\label{def:srht}
Sea $n$ una potencia de $2$. La SRHT (\emph{subsampled randomized Hadamard transform}) es
\[ S=\sqrt{\frac{n}{m}}\;R\,H\,D, \]
donde
\begin{itemize}
  \item $D\in\R^{n\times n}$ es diagonal con signos $\pm1$ independientes y uniformes;
  \item $H\in\R^{n\times n}$ es la matriz de Walsh--Hadamard normalizada, definida por
        $H_1=(1)$, $H_{2k}=\frac{1}{\sqrt2}\begin{pmatrix}H_k&H_k\\H_k&-H_k\end{pmatrix}$; es ortogonal y todas sus entradas valen $\pm1/\sqrt n$;
  \item $R\in\R^{m\times n}$ selecciona $m$ filas al azar, uniformemente.
\end{itemize}
\end{definicion}

\paragraph{Cómo funciona.} El muestreo uniforme de filas es el \emph{sketch} más barato posible, pero fracasa si
unas pocas filas concentran la información (sección~\ref{sec:muestreo}). La SRHT lo arregla
\emph{mezclando antes de muestrear}: la rotación aleatoria $HD$ reparte la «masa» de cada vector entre todas
las coordenadas. Es una manifestación del principio de incertidumbre: un vector concentrado en pocas
coordenadas tiene transformada de Hadamard dispersa; los signos aleatorios $D$ impiden que un vector esté
adversarialmente alineado con $H$. En el lenguaje del capítulo~\ref{cap:embed}, $HD$ \emph{uniformiza las
puntuaciones de palanca}: con probabilidad al menos $1-p$,
\begin{equation}\label{eq:srht-aplana}
  \max_{1\le i\le n}\nrm{e_i\T HDU}_2\;\le\;\sqrt{\frac dn}+\sqrt{\frac{8\log(n/p)}{n}} ,
\end{equation}
mientras que el valor medio de $\nrm{e_i\T HDU}_2^2$ es $d/n$ (Tropp, 2011). Tras la mezcla, todas las filas son
casi igual de importantes y el muestreo uniforme es seguro.

\begin{ejemplo}[SRHT con $n=4$]
\[ H_4=\frac12\begin{pmatrix}1&1&1&1\\1&-1&1&-1\\1&1&-1&-1\\1&-1&-1&1\end{pmatrix}. \]
Para el vector $x=e_1=(1,0,0,0)\T$, concentrado en una coordenada, $HDx=\pm\frac12(1,1,1,1)\T$: perfectamente
disperso. Muestrear $m=2$ coordenadas de $x$ daría norma $0$ con probabilidad $\tfrac12$; muestrear $2$
coordenadas de $HDx$ y reescalar por $\sqrt{n/m}=\sqrt2$ da siempre norma exactamente $1$.
\end{ejemplo}

\paragraph{Coste.} Multiplicar por $D$ cuesta $\cO(nd)$; aplicar $H$ a las $d$ columnas con la transformada
rápida de Walsh--Hadamard cuesta $\cO(nd\log n)$ (y $\cO(nd\log m)$ con una variante podada que calcula solo
las $m$ coordenadas necesarias). No se almacena ninguna matriz. Si $n$ no es potencia de $2$ se rellena $A$
con filas de ceros. Otras transformadas ortogonales con entradas de módulo $1/\sqrt n$ (Fourier, coseno)
sirven igual; la familia se conoce como \emph{sistemas ortogonales aleatorizados} (ROS).

\section{\emph{Sketches} dispersos: CountSketch, SJLT y OSNAP}

\begin{definicion}[CountSketch y SJLT]\label{def:countsketch}
En una SJLT (\emph{sparse Johnson--Lindenstrauss transform}) con $s$ no nulos por columna, cada columna de
$S\in\R^{m\times n}$ tiene exactamente $s$ entradas no nulas, situadas en filas elegidas al azar y con valores
$\pm1/\sqrt s$ con signos independientes. El caso $s=1$ es \textbf{CountSketch}: $S_{h(i),i}=\sigma_i$, con
$h:\{1,\dots,n\}\to\{1,\dots,m\}$ una función de dispersión (\emph{hash}) aleatoria y $\sigma_i=\pm1$.
\end{definicion}

\paragraph{Cómo funciona.} Con $s=1$, calcular $SA$ consiste en recorrer las filas de $A$ una sola vez: la fila
$i$ se multiplica por el signo $\sigma_i$ y se \emph{suma} a la fila $h(i)$ de $SA$ («cubeta» $h(i)$).
\[ \underbrace{\begin{pmatrix}0&-1&0&0&+1&0\\+1&0&0&-1&0&0\\0&0&+1&0&0&+1\end{pmatrix}}_{S\ (m=3,\ n=6)}
   \begin{pmatrix}a_1\T\\a_2\T\\a_3\T\\a_4\T\\a_5\T\\a_6\T\end{pmatrix}
   =\begin{pmatrix}-a_2\T+a_5\T\\ a_1\T-a_4\T\\ a_3\T+a_6\T\end{pmatrix}. \]
Los signos aleatorios hacen que los términos cruzados se cancelen en media:
$\E[(SA)\T(SA)]=\sum_{i,j}\E[\sigma_i\sigma_j\mathbf 1\{h(i)=h(j)\}]\,a_ia_j\T=\sum_ia_ia_i\T=A\T A$.

\paragraph{Coste.} $\cO(\nnz(A))$ para $s=1$ y $\cO(s\cdot\nnz(A))$ en general: \emph{tiempo proporcional a la
dispersidad de los datos}, lo mejor posible. Es un único pase por los datos y admite datos en flujo.

\paragraph{El precio.} Dos filas «importantes» que caen en la misma cubeta se mezclan de forma irreversible.
Si el subespacio tiene $d$ direcciones importantes, evitar colisiones entre ellas exige, por la paradoja del
cumpleaños, $m\gtrsim d^2$ cubetas. Esta dependencia cuadrática es inevitable para $s=1$; tomando
$s\asymp\log d$ no nulos por columna (construcción OSNAP de Nelson y Nguyên) se recupera $m\asymp d\log d$.
El experimento de la sección~\ref{sec:exp-embed} muestra el fenómeno con nitidez.

\section{Muestreo de filas}\label{sec:muestreo}

\begin{definicion}[\emph{Sketch} de muestreo]
Dada una distribución $(p_1,\dots,p_n)$ sobre las filas, cada fila de $S$ es, independientemente,
$e_i\T/\sqrt{m\,p_i}$ con probabilidad $p_i$. Así, $SA$ consta de $m$ filas de $A$ reescaladas.
\end{definicion}

El reescalado garantiza \eqref{eq:isotropia}. Las tres elecciones habituales son:
\begin{itemize}
  \item \textbf{Uniforme}, $p_i=1/n$. Coste nulo. Funciona si y solo si $A$ es \emph{incoherente} (ninguna fila
        es mucho más importante que las demás).
  \item \textbf{Por puntuaciones de palanca}, $p_i=\ell_i/d$ con $\ell_i=a_i\T(A\T A)^{-1}a_i$. Funciona siempre
        con $m\asymp d\log d/\eps^2$, pero calcular $\ell_i$ exactamente cuesta tanto como resolver el problema;
        se usan aproximaciones.
  \item \textbf{Por norma de filas}, $p_i\propto\nrm{a_i}_2^2$. Barato; garantías más débiles (en norma de Frobenius).
\end{itemize}
A diferencia de los anteriores, los dos últimos \emph{no son ajenos a los datos}: $S$ depende de $A$.

\begin{observacion}[Newton submuestreado]\label{obs:subsampled}
Cuando $f(x)=\frac1n\sum_i f_i(x)$, aproximar el hessiano con un subconjunto aleatorio de sumandos es
exactamente aplicar un \emph{sketch} de muestreo uniforme a la raíz del hessiano. Los métodos de Newton
submuestreado (Erdogdu y Montanari; Roosta-Khorasani y Mahoney) son, por tanto, casos particulares de Newton
Sketch. En optimización estocástica por escenarios, \emph{muestrear escenarios} es muestrear filas
(capítulo~\ref{cap:estoc}).
\end{observacion}

\section{Comparación}

\begin{table}[ht]
\centering\small
\caption{Matrices de \emph{sketch} principales. La columna $m$ indica el número de filas suficiente para un
encaje de subespacio con distorsión $\eps$ y probabilidad de fallo $p$ (teoremas del capítulo~\ref{cap:embed}).}
\label{tab:sketches}
\begin{tabular}{@{}llll@{}}
\toprule
\emph{Sketch} & Tamaño $m$ suficiente & Coste de $SA$ & ¿Ajeno a los datos? \\
\midrule
Gaussiano & $\dfrac{d+\log(1/p)}{\eps^2}$ & $mnd$ & sí \\[8pt]
SRHT & $\dfrac{(d+\log(n/p))\log(d/p)}{\eps^2}$ & $nd\log n$ & sí \\[8pt]
CountSketch ($s=1$) & $\dfrac{d^2}{\eps^2\,p}$ & $\nnz(A)$ & sí \\[8pt]
OSNAP ($s\asymp\log(d/p)/\eps$) & $\dfrac{d\log(d/p)}{\eps^2}$ & $s\cdot\nnz(A)$ & sí \\[8pt]
Muestreo por palanca & $\dfrac{d\log(d/p)}{\eps^2}$ & $\nnz(A)$ + cálculo de $\ell_i$ & no \\[8pt]
Muestreo uniforme & $\dfrac{\mu(A)\,d\log(d/p)}{\eps^2}$ & $md$ & sí \\
\bottomrule
\end{tabular}
\end{table}

En la tabla~\ref{tab:sketches} se omiten constantes universales y $\mu(A)\in[1,n/d]$ es la coherencia
(definición~\ref{def:palanca}). La elección práctica suele ser: SRHT para datos densos, CountSketch o SJLT
para datos dispersos o cuando el tiempo de aplicación domina, y muestreo cuando solo se puede acceder a filas.

\begin{resumen}
\begin{itemize}
  \item Un \emph{sketch} $S$ ($m\times n$, $\E[S\T S]=I$) comprime $A$ a $SA$ conservando normas en media.
  \item Gaussiano: el mejor $m$ ($\asymp d/\eps^2$) y el peor coste ($mnd$).
  \item SRHT: mezcla ($HD$) y luego muestrea; $m\asymp d\log d/\eps^2$, coste $nd\log n$.
  \item CountSketch: suma con signos en cubetas; coste $\nnz(A)$, pero $m\asymp d^2/\eps^2$.
  \item Muestreo: trivial de aplicar; necesita incoherencia o probabilidades adaptadas (palanca).
\end{itemize}
\end{resumen}

\chapter{Teoría de subespacios encajados}\label{cap:embed}

El término inglés es \emph{subspace embedding}; en español se encuentra como «encaje de subespacio»,
«inmersión de subespacio» o, coloquialmente, «subembedding». La idea: una matriz $S$ con muy pocas filas
que, aplicada a \emph{todos} los vectores de un subespacio de dimensión $d$, conserva sus normas salvo un
factor $1\pm\eps$. Es la herramienta que convierte «$S$ conserva normas en media» en «$SA$ puede sustituir a
$A$ en un algoritmo».

\section{El punto de partida: Johnson--Lindenstrauss}

\begin{lema}[Concentración para un vector fijo]\label{lem:jl-dist}
Sea $S\in\R^{m\times n}$ gaussiana (entradas $N(0,1/m)$ independientes) y $x\in\R^n$ fijo. Para $\eps\in(0,1)$,
\[ \Prob\Bigl(\,\bigl|\nrm{Sx}_2^2-\nrm{x}_2^2\bigr|>\eps\nrm{x}_2^2\Bigr)\;\le\;2\exp\bigl(-m\eps^2/12\bigr). \]
\end{lema}

\begin{proof}
Por homogeneidad, $\nrm{x}_2=1$. Entonces $Sx\sim N(0,I_m/m)$ y $m\nrm{Sx}_2^2\sim\chi^2_m$. Las cotas de Chernoff
para la $\chi^2$ (apéndice~\ref{ap:conc}) dan
$\Prob(\chi^2_m\ge(1+\eps)m)\le\exp(-m(\eps^2/4-\eps^3/6))\le e^{-m\eps^2/12}$ y $\Prob(\chi^2_m\le(1-\eps)m)\le e^{-m\eps^2/4}$.
\end{proof}

\begin{teorema}[Lema de Johnson--Lindenstrauss]\label{teo:jl}
Sean $x_1,\dots,x_N\in\R^n$ y $\eps,p\in(0,1)$. Si $m\ge12\,\eps^{-2}\log(N^2/p)$ y $S$ es gaussiana, entonces
con probabilidad al menos $1-p$,
\[ (1-\eps)\nrm{x_i-x_j}_2^2\;\le\;\nrm{Sx_i-Sx_j}_2^2\;\le\;(1+\eps)\nrm{x_i-x_j}_2^2\qquad\text{para todo } i,j. \]
\end{teorema}

\begin{proof}
Se aplica el lema~\ref{lem:jl-dist} a cada uno de los $\binom N2<N^2/2$ vectores $x_i-x_j$ y la cota de la unión:
la probabilidad de fallo es a lo sumo $N^2e^{-m\eps^2/12}\le p$.
\end{proof}

Lo notable es que $m$ depende de $\log N$ y \emph{no} de la dimensión ambiente $n$. Pero el lema habla de un
conjunto \emph{finito}; un subespacio tiene infinitos vectores. El paso de uno a otro es el contenido de este capítulo.

\section{Definición y formas equivalentes}

\begin{definicion}[Encaje de subespacio]\label{def:embed}
Sea $A\in\R^{n\times d}$ y $\eps\in(0,1)$. Una matriz $S\in\R^{m\times n}$ es un \textbf{encaje de subespacio
con distorsión $\eps$} para $A$ (para $\col(A)$) si
\begin{equation}\label{eq:embed}
  (1-\eps)\nrm{Ax}_2^2\;\le\;\nrm{SAx}_2^2\;\le\;(1+\eps)\nrm{Ax}_2^2\qquad\text{para todo } x\in\R^d .
\end{equation}
Una distribución sobre matrices $S$ es un \textbf{encaje de subespacio ajeno a los datos} (\emph{oblivious
subspace embedding}, OSE) con parámetros $(\eps,p,d)$ si, para \emph{cada} $A$ con $d$ columnas fijada de
antemano, $S$ cumple \eqref{eq:embed} con probabilidad al menos $1-p$.
\end{definicion}

\begin{proposicion}[Formas equivalentes]\label{prop:embed-equiv}
Sea $A$ de rango $d$ y $U\in\R^{n\times d}$ una base ortonormal de $\col(A)$. Son equivalentes:
\begin{enumerate}[label=(\roman*)]
  \item $S$ es un encaje con distorsión $\eps$ para $A$, es decir, \eqref{eq:embed}.
  \item $\nrm{U\T S\T SU-I_d}_2\le\eps$.
  \item Todos los valores singulares de $SU$ están en $[\sqrt{1-\eps},\sqrt{1+\eps}]$.
  \item \textbf{Forma de Loewner:}\quad $(1-\eps)\,A\T A\;\preceq\;A\T S\T SA\;\preceq\;(1+\eps)\,A\T A$.
\end{enumerate}
\end{proposicion}

\begin{proof}
Todo vector de $\col(A)$ es $Uz$ con $\nrm{Uz}_2=\nrm{z}_2$, así que (i) dice
\[ \abs{z\T(U\T S\T SU-I)z}\le\eps\nrm{z}_2^2\qquad\text{para todo }z, \]
que para una matriz simétrica es (ii); y (ii) dice que
los valores propios de $U\T S\T SU$, que son los cuadrados de los valores singulares de $SU$, están en
$[1-\eps,1+\eps]$: es (iii). Por último, (i) escrita como $(1-\eps)x\T A\T Ax\le x\T A\T S\T SAx\le(1+\eps)x\T A\T Ax$ es (iv).
\end{proof}

\begin{ideaclave}[La forma (iv) es el puente con la optimización]
Si $A=B(x)$ es la raíz cuadrada del hessiano, (iv) dice
\[ (1-\eps)\,\nabla^2 f(x)\;\preceq\;B\T S\T SB\;\preceq\;(1+\eps)\,\nabla^2 f(x): \]
el hessiano con \emph{sketch} aproxima al exacto \emph{espectralmente y en sentido relativo}, que es
exactamente el tipo de error que el método de Newton tolera (capítulos~\ref{cap:ns} y~\ref{cap:ans}).
\end{ideaclave}

\section{De un vector a un subespacio: el argumento de la \texorpdfstring{$\eps$}{eps}-red}

\begin{lema}[Basta comprobar en una red]\label{lem:red}
Sea $M\in\R^{d\times d}$ simétrica y $\cN$ una $\tfrac14$-red de la esfera unidad $\mathbb S^{d-1}$ (todo punto
de la esfera dista a lo sumo $\tfrac14$ de algún punto de $\cN$). Entonces $\nrm{M}_2\le2\max_{y\in\cN}\abs{y\T My}$.
Existe una tal red con $\abs{\cN}\le9^d$.
\end{lema}

\begin{proof}
Sea $x$ unitario con $\abs{x\T Mx}=\nrm{M}_2$ e $y\in\cN$ con $\nrm{x-y}_2\le\tfrac14$. Entonces
$\abs{x\T Mx-y\T My}=\abs{x\T M(x-y)+(x-y)\T My}\le2\cdot\tfrac14\nrm{M}_2$, luego $\abs{y\T My}\ge\tfrac12\nrm{M}_2$.
La cota de cardinalidad es el argumento volumétrico habitual: $(1+2/\rho)^d$ con $\rho=\tfrac14$.
\end{proof}

\begin{teorema}[El \emph{sketch} gaussiano es un encaje de subespacio]\label{teo:gauss-ose}
Sea $S\in\R^{m\times n}$ gaussiana y $\eps,p\in(0,1)$. Si
\[ m\;\ge\;\frac{48}{\eps^2}\Bigl(d\log9+\log\frac2p\Bigr), \]
entonces, para cada $A\in\R^{n\times d}$ fija, $S$ es un encaje de subespacio con distorsión $\eps$ para $A$ con
probabilidad al menos $1-p$. Es decir, basta $m=\cO\bigl((d+\log(1/p))/\eps^2\bigr)$.
\end{teorema}

\begin{proof}
Sea $U$ una base ortonormal de $\col(A)$ y $M=U\T S\T SU-I_d$. Para $y$ unitario,
$y\T My=\nrm{SUy}_2^2-1$, y $Uy$ es un vector unitario \emph{fijo} de $\R^n$: por el lema~\ref{lem:jl-dist} con
$\eps/2$, $\Prob(\abs{y\T My}>\eps/2)\le2e^{-m\eps^2/48}$. Por la cota de la unión sobre la red del
lema~\ref{lem:red}, con probabilidad al menos $1-2\cdot9^de^{-m\eps^2/48}\ge1-p$ se tiene
$\max_{y\in\cN}\abs{y\T My}\le\eps/2$, y entonces $\nrm{M}_2\le\eps$, que es la forma (ii) de la proposición~\ref{prop:embed-equiv}.
\end{proof}

El mecanismo es general: \emph{concentración para un vector} $+$ \emph{red de tamaño exponencial en $d$}
$\Rightarrow$ $m$ lineal en $d$. La dependencia $d/\eps^2$ es óptima: cualquier OSE necesita
$m=\Omega(d/\eps^2)$ filas (Nelson y Nguyên, 2014).

\section{Puntuaciones de palanca, coherencia y muestreo}

\begin{definicion}[Puntuaciones de palanca y coherencia]\label{def:palanca}
Sea $A\in\R^{n\times d}$ de rango $d$ y $U$ una base ortonormal de $\col(A)$. La \textbf{puntuación de palanca}
(\emph{leverage score}) de la fila $i$ es
\[ \ell_i=\nrm{e_i\T U}_2^2=a_i\T(A\T A)^{-1}a_i . \]
La \textbf{coherencia} de $A$ es $\mu(A)=\frac nd\max_i\ell_i$.
\end{definicion}

Propiedades inmediatas: $0\le\ell_i\le1$; $\sum_i\ell_i=\nrm{U}_F^2=d$; no dependen de la base elegida; y
$1\le\mu(A)\le n/d$. Una fila con $\ell_i\approx1$ es \emph{imprescindible}: hay una dirección del espacio columna
que solo ella «ve». Una matriz con $\mu(A)\approx1$ es \emph{incoherente}: todas las filas importan por igual.

\begin{teorema}[Muestreo por palanca y muestreo uniforme]\label{teo:muestreo}
Sea $S$ un \emph{sketch} de muestreo con probabilidades $p_i$ y $\eps,p\in(0,1)$.
\begin{enumerate}[label=(\alph*)]
  \item Si $p_i=\ell_i/d$ y $m\ge3\,d\,\eps^{-2}\log(2d/p)$, entonces $S$ es un encaje con distorsión $\eps$ para $A$ con probabilidad $\ge1-p$.
  \item Si $p_i=1/n$ y $m\ge3\,\mu(A)\,d\,\eps^{-2}\log(2d/p)$, la misma conclusión.
\end{enumerate}
\end{teorema}

\begin{proof}
Sea $u_i=U\T e_i\in\R^d$ (fila $i$ de $U$). Entonces $U\T S\T SU=\sum_{k=1}^mX_k$ con
$X_k=\frac{1}{m\,p_{i_k}}u_{i_k}u_{i_k}\T$ independientes y semidefinidas positivas, y
$\E X_k=\frac1m\sum_iu_iu_i\T=\frac1mI_d$, de modo que $\E\sum_kX_k=I_d$. Además
$\lambda_{\max}(X_k)=\ell_{i_k}/(m\,p_{i_k})\le R$, con $R=d/m$ en el caso (a) y $R=\mu(A)d/m$ en el caso (b). La
desigualdad de Chernoff matricial (teorema~\ref{teo:chernoff-mat}) con $\mu_{\min}=\mu_{\max}=1$ da
$\Prob(\nrm{\sum_kX_k-I}_2>\eps)\le2d\exp(-\eps^2/(3R))$, que es $\le p$ bajo las hipótesis.
\end{proof}

El factor $\log d$ es aquí inevitable (problema del coleccionista de cupones: hay que «ver» las $d$
direcciones). El teorema también explica la SRHT: por \eqref{eq:srht-aplana}, tras la mezcla $HD$ la
coherencia es $\mu(HDA)\le\bigl(1+\sqrt{8\log(n/p)/d}\bigr)^2$, y el apartado (b) da el tamaño de la tabla~\ref{tab:sketches}.

\section{Teoremas principales para \emph{sketches} rápidos}

\begin{teorema}[SRHT {\cite{Tropp11,BG13}}]\label{teo:srht}
Sea $S$ una SRHT y $\eps,p\in(0,1)$. Si
\[ m\;\gtrsim\;\frac{1}{\eps^2}\Bigl(\sqrt d+\sqrt{8\log(n/p)}\Bigr)^{2}\log(d/p), \]
entonces $S$ es un encaje con distorsión $\eps$ para cada $A\in\R^{n\times d}$ fija con probabilidad $\ge1-p$.
\end{teorema}

\begin{teorema}[CountSketch {\cite{CW13,NN13,MM13}}]\label{teo:count}
Sea $S$ un CountSketch y $\eps,p\in(0,1)$. Si $m\ge2d^2/(\eps^2p)$, entonces $S$ es un encaje con distorsión
$\eps$ para cada $A\in\R^{n\times d}$ fija con probabilidad $\ge1-p$.
\end{teorema}

\begin{proof}
Método del segundo momento. Con $u_i$ las filas de $U$,
$(U\T S\T SU-I)_{ab}=\sum_{i\ne j}\sigma_i\sigma_j\mathbf 1\{h(i)=h(j)\}U_{ia}U_{jb}$. Al tomar la esperanza del
cuadrado solo sobreviven los pares $\{i,j\}=\{i',j'\}$, cada uno con probabilidad de colisión $1/m$:
\[ \E\nrm{U\T S\T SU-I}_F^2=\frac1m\sum_{i\ne j}\bigl(\ell_i\ell_j+(u_i\T u_j)^2\bigr)\le\frac1m\bigl(d^2+\nrm{UU\T}_F^2\bigr)=\frac{d^2+d}{m}\le\frac{2d^2}{m}. \]
Por Markov, $\Prob(\nrm{U\T S\T SU-I}_2>\eps)\le\Prob(\nrm{\cdot}_F^2>\eps^2)\le2d^2/(m\eps^2)\le p$.
\end{proof}

La demostración solo usa independencia de a pares en $h$ y de a cuatro en los signos. La dependencia $d^2$ no
es un defecto del análisis: para $s=1$ no nulo por columna hacen falta $m=\Omega(d^2)$ filas (Nelson y Nguyên).

\begin{teorema}[OSNAP {\cite{NN13,Cohen16}}]\label{teo:osnap}
Una SJLT con $m\gtrsim d\log(d/p)/\eps^2$ filas y $s\gtrsim\log(d/p)/\eps$ no nulos por columna es un encaje de
subespacio con distorsión $\eps$ para cada $A\in\R^{n\times d}$ fija con probabilidad $\ge1-p$.
\end{teorema}

\section{Encajes regularizados y dimensión efectiva}\label{sec:dim-efectiva}

En muchos problemas el hessiano es $A\T A+\mu I$ (regresión \emph{ridge}) o, más en general,
$B\T B+\nabla^2 g$ con $\nabla^2 g\succeq\mu I$. Solo se aplica el \emph{sketch} a $B$, y lo que se necesita es
\begin{equation}\label{eq:embed-reg}
  (1-\eps)\,(A\T A+\mu I)\;\preceq\;A\T S\T SA+\mu I\;\preceq\;(1+\eps)\,(A\T A+\mu I).
\end{equation}
Esto es \emph{más débil} que \eqref{eq:embed}: en las direcciones donde $A\T A\ll\mu I$, el término $\mu I$ ya
domina y da igual lo que haga el \emph{sketch}.

\begin{definicion}[Dimensión efectiva]\label{def:dim-efectiva}
La \textbf{dimensión efectiva} de $A$ con regularización $\mu>0$ es
\[ \dmu=\tr\bigl(A\T A\,(A\T A+\mu I)^{-1}\bigr)=\sum_{j=1}^d\frac{\sigma_j^2}{\sigma_j^2+\mu}, \]
con $\sigma_1\ge\dots\ge\sigma_d$ los valores singulares de $A$.
\end{definicion}

Cada valor singular aporta un número entre $0$ y $1$: casi $1$ si $\sigma_j^2\gg\mu$, casi $0$ si
$\sigma_j^2\ll\mu$. Así, $\dmu$ \emph{cuenta las direcciones en las que los datos pesan más que la
regularización}. Siempre $\dmu\le\rank(A)\le d$ y $\dmu\le\nrm{A}_F^2/\mu$.

\begin{ejemplo}
Si $\sigma_j^2=C\,j^{-2\beta}$ con $\beta>\tfrac12$, entonces $\dmu\asymp(C/\mu)^{1/(2\beta)}$, que puede ser
mucho menor que $d$ e incluso no depender de $d$. En el experimento de la sección~\ref{sec:exp-ans},
$d=256$ y $\dmu\approx26$.
\end{ejemplo}

\begin{proposicion}[Por qué manda la dimensión efectiva]\label{prop:embed-reg}
Sea $H=A\T A+\mu I$ y $Z=AH^{-1/2}\in\R^{n\times d}$. Entonces \eqref{eq:embed-reg} equivale a
\[ \nrm{Z\T(S\T S-I_n)Z}_2\le\eps,\qquad\text{y}\qquad\nrm{Z}_2^2=\frac{\sigma_1^2}{\sigma_1^2+\mu}\le1,\qquad\nrm{Z}_F^2=\dmu . \]
\end{proposicion}

\begin{proof}
Multiplicando \eqref{eq:embed-reg} a ambos lados por $H^{-1/2}$ se obtiene
$-\eps I\preceq H^{-1/2}(A\T S\T SA-A\T A)H^{-1/2}\preceq\eps I$, y el término central es $Z\T(S\T S-I)Z$. Las dos
normas se leen en la descomposición en valores singulares de $A$.
\end{proof}

La matriz $Z$ tiene norma espectral a lo sumo $1$ y norma de Frobenius al cuadrado $\dmu$: su \emph{rango
estable} es del orden de $\dmu$. Los resultados de producto matricial aproximado en términos del rango
estable (Cohen, Nelson y Woodruff, 2016) y de regresión regularizada (Avron, Clarkson y Woodruff, 2017)
muestran que entonces \textbf{basta un \emph{sketch} de tamaño proporcional a $\dmu$ en lugar de $d$}:
$m\asymp\dmu\log\dmu/\eps^2$ para la SRHT y $m\asymp\dmu^2/\eps^2$ para CountSketch (el enunciado preciso
que usa Newton Sketch adaptativo es el lema~\ref{lem:lwp1}). El mismo cálculo con muestreo lleva a las
\emph{puntuaciones de palanca regularizadas} $\ell_i^{\mu}=a_i\T(A\T A+\mu I)^{-1}a_i$, cuya suma es $\dmu$.



\begin{resumen}
\begin{itemize}
  \item Encaje de subespacio: $\nrm{SAx}^2\in(1\pm\eps)\nrm{Ax}^2$ para todo $x$; equivale a $(1-\eps)A\T A\preceq A\T S\T SA\preceq(1+\eps)A\T A$.
  \item Demostración tipo: concentración para un vector $+$ $\eps$-red $\Rightarrow$ $m\asymp d/\eps^2$ (gaussiano).
  \item Muestreo: Chernoff matricial; el tamaño lo fija $\max_i\ell_i/p_i$. La SRHT uniformiza las palancas.
  \item CountSketch: segundo momento; $m\asymp d^2/\eps^2$ por colisiones.
  \item Con regularización $\mu$, el tamaño necesario lo fija la dimensión efectiva $\dmu=\sum_j\sigma_j^2/(\sigma_j^2+\mu)\le d$.
\end{itemize}
\end{resumen}

\chapter{Aplicación a mínimos cuadrados}\label{cap:ls}

El problema de mínimos cuadrados sobredeterminado,
\begin{equation}\label{eq:ls}
  x^\star=\argmin_{x\in\R^d}\;\tfrac12\nrm{Ax-b}_2^2,\qquad A\in\R^{n\times d},\ n\gg d,
\end{equation}
es el laboratorio natural del \emph{sketching}: es el caso cuadrático de \eqref{eq:problema-general}, con hessiano
constante $A\T A$. Las tres estrategias que aparecen aquí (\emph{resolver el problema comprimido},
\emph{precondicionar} e \emph{iterar con el hessiano comprimido}) reaparecen después en optimización general, y
la tercera \emph{es} Newton Sketch.

\section{Métodos clásicos y su coste}

\begin{itemize}
  \item \textbf{Directos} (QR de Householder, SVD): $\cO(nd^2)$ operaciones, numéricamente estables.
  \item \textbf{Ecuaciones normales} $A\T Ax=A\T b$ con Cholesky: $\cO(nd^2)$, pero el número de condición es
        $\kappa(A)^2$ y se pierde el doble de cifras significativas.
  \item \textbf{Iterativos} (gradiente conjugado sobre las ecuaciones normales, LSQR): $\cO(nd)$ por iteración y del
        orden de $\kappa(A)\log(1/\delta)$ iteraciones.
\end{itemize}
Los directos pagan $nd^2$; los iterativos pagan $\kappa(A)$. El \emph{sketching} permite no pagar ninguno de los dos.

\section{Primera estrategia: \emph{sketch-and-solve}}

Se resuelve el problema comprimido:
\begin{equation}\label{eq:sas}
  \tilde x=\argmin_{x\in\R^d}\;\tfrac12\nrm{S(Ax-b)}_2^2 .
\end{equation}
Coste: el de formar $SA$ y $Sb$ más $\cO(md^2)$.

\begin{teorema}[Garantía de \emph{sketch-and-solve}]\label{teo:sas}
Si $S$ es un encaje de subespacio con distorsión $\eps$ para la matriz ampliada $[A\ \ b]\in\R^{n\times(d+1)}$, entonces
\[ \nrm{A\tilde x-b}_2^2\;\le\;\frac{1+\eps}{1-\eps}\,\nrm{Ax^\star-b}_2^2 . \]
\end{teorema}

\begin{proof}
Para todo $x$, el vector $Ax-b$ está en $\col([A\ b])$. Entonces
\[ \nrm{A\tilde x-b}_2^2\le\frac{1}{1-\eps}\nrm{S(A\tilde x-b)}_2^2\le\frac{1}{1-\eps}\nrm{S(Ax^\star-b)}_2^2\le\frac{1+\eps}{1-\eps}\nrm{Ax^\star-b}_2^2, \]
donde la desigualdad central usa que $\tilde x$ minimiza el problema comprimido.
\end{proof}

Este argumento exige $m\asymp d/\eps^2$. Un análisis más fino (Sarlós; Drineas, Mahoney, Muthukrishnan y
Sarlós) muestra que para el \emph{residuo} basta $m\asymp d/\eps$. En el caso gaussiano el cálculo es exacto:

\begin{proposicion}[Exceso de residuo con \emph{sketch} gaussiano]\label{prop:sas-gauss}
Si $S$ es gaussiana con $m>d+1$ y $A$ tiene rango $d$,
\[ \E\,\nrm{A\tilde x-b}_2^2=\Bigl(1+\frac{d}{m-d-1}\Bigr)\nrm{Ax^\star-b}_2^2 . \]
\end{proposicion}

\begin{proof}
Sea $r^\star=b-Ax^\star\perp\col(A)$ y $A=UT$ con $U$ ortonormal. Las ecuaciones normales del problema
comprimido dan $A(\tilde x-x^\star)=U(U\T S\T SU)^{-1}U\T S\T Sr^\star$. Por invariancia rotacional, $G_1=SU$ y
$g=Sr^\star/\nrm{r^\star}$ son gaussianos independientes con entradas de varianza $1/m$. Condicionando en $G_1$,
$\E\nrm{(G_1\T G_1)^{-1}G_1\T g}_2^2=\frac1m\tr\bigl((G_1\T G_1)^{-1}\bigr)$, y como $G_1\T G_1$ es una matriz de Wishart
con $\E[(G_1\T G_1)^{-1}]=\frac{m}{m-d-1}I_d$, resulta $\E\nrm{A(\tilde x-x^\star)}_2^2=\frac{d}{m-d-1}\nrm{r^\star}_2^2$.
Por Pitágoras, $\nrm{A\tilde x-b}^2=\nrm{r^\star}^2+\nrm{A(\tilde x-x^\star)}^2$.
\end{proof}

\begin{cuidado}[El límite de \emph{sketch-and-solve}]
El error relativo en el residuo es $\approx d/m$. Para tres cifras ($10^{-3}$) haría falta $m\approx1000\,d$;
para diez cifras, $m\approx10^{10}d\gg n$. \emph{Sketch-and-solve} solo da soluciones de baja precisión. Para alta
precisión hay que \emph{iterar}, usando el \emph{sketch} solo para acelerar cada iteración.
\end{cuidado}

\section{Segunda estrategia: \emph{sketch-and-precondition}}

Idea (Rokhlin y Tygert; Blendenpik; LSRN): usar $SA$ no para resolver, sino para construir un
precondicionador. Se calcula la factorización $SA=QR$ ($R\in\R^{d\times d}$) y se resuelve con LSQR el
problema equivalente $\min_y\nrm{(AR^{-1})y-b}_2$, $x=R^{-1}y$.

\begin{teorema}[Calidad del precondicionador]\label{teo:precond}
Si $S$ es un encaje de subespacio con distorsión $\eps$ para $A$ y $SA=QR$, entonces
\[ \kappa(AR^{-1})\;\le\;\sqrt{\frac{1+\eps}{1-\eps}} . \]
\end{teorema}

\begin{proof}
Para $y\in\R^d$, $\nrm{SAR^{-1}y}_2=\nrm{Qy}_2=\nrm{y}_2$. Por \eqref{eq:embed} aplicada a $x=R^{-1}y$,
$\frac{1}{1+\eps}\nrm{y}_2^2\le\nrm{AR^{-1}y}_2^2\le\frac{1}{1-\eps}\nrm{y}_2^2$.
\end{proof}

De hecho los valores singulares de $AR^{-1}$ son exactamente los inversos de los de $SU$; para un \emph{sketch}
gaussiano, por \eqref{eq:gauss-sv}, $\kappa(AR^{-1})\approx\frac{1+\sqrt{d/m}}{1-\sqrt{d/m}}$: con $m=4d$ es $3$, \textbf{sea
cual sea $\kappa(A)$}. LSQR reduce el error en un factor $\frac{\kappa-1}{\kappa+1}\approx\sqrt{d/m}$ por iteración,
así que bastan $\cO(\log(1/\delta))$ iteraciones. Coste total con SRHT y $m\asymp d$:
\[ \underbrace{\cO(nd\log d)}_{\text{sketch}}+\underbrace{\cO(d^3)}_{\text{QR de }SA}+\underbrace{\cO(nd\log(1/\delta))}_{\text{LSQR}}, \]
frente a $\cO(nd^2)$ de un método directo. Aquí basta una distorsión \emph{constante} (por ejemplo
$\eps=\tfrac12$): la precisión final la da la iteración, no el \emph{sketch}.

\section{Tercera estrategia: \emph{sketch} iterativo del hessiano (IHS)}\label{sec:ihs}

Pilanci y Wainwright (2016) observaron que el defecto de \emph{sketch-and-solve} es comprimir también $b$. Si
se escribe el paso de Newton exacto para \eqref{eq:ls} desde $x_t$,
\[ x_{t+1}=\argmin_x\Bigl\{\tfrac12\nrm{A(x-x_t)}_2^2-\ip{A\T(b-Ax_t)}{x-x_t}\Bigr\}, \]
y se aplica el \emph{sketch} \textbf{solo al término cuadrático} (el hessiano), manteniendo el gradiente exacto, se obtiene
\begin{equation}\label{eq:ihs}
  x_{t+1}=\argmin_x\Bigl\{\tfrac12\nrm{S_tA(x-x_t)}_2^2-\ip{A\T(b-Ax_t)}{x-x_t}\Bigr\}
  =x_t+(A\T S_t\T S_tA)^{-1}A\T(b-Ax_t).
\end{equation}

\begin{teorema}[Convergencia lineal del IHS]\label{teo:ihs}
Si $S_t$ es un encaje de subespacio con distorsión $\eps<\tfrac12$ para $A$, entonces
\[ \nrm{A(x_{t+1}-x^\star)}_2\;\le\;\frac{\eps}{1-\eps}\,\nrm{A(x_t-x^\star)}_2 . \]
\end{teorema}

\begin{proof}
Como $A\T(b-Ax_t)=-A\T A(x_t-x^\star)$, \eqref{eq:ihs} da
\[ x_{t+1}-x^\star=\bigl(I-(A\T S_t\T S_tA)^{-1}A\T A\bigr)(x_t-x^\star). \]
Con $A=UT$ y $W=U\T S_t\T S_tU$, y llamando $e_t=T(x_t-x^\star)$ (de modo que $\nrm{e_t}_2=\nrm{A(x_t-x^\star)}_2$), queda
$e_{t+1}=(I-W^{-1})e_t$. Los valores propios de $W$ están en $[1-\eps,1+\eps]$, luego los de $I-W^{-1}$ están en
$[-\frac{\eps}{1-\eps},\frac{\eps}{1+\eps}]$.
\end{proof}

\begin{ideaclave}[Por qué el IHS llega a cualquier precisión y \emph{sketch-and-solve} no]
El punto fijo de \eqref{eq:ihs} es $A\T(b-Ax)=0$: la solución \emph{exacta}. El \emph{sketch} solo afecta a la
\emph{velocidad} (el factor $\frac{\eps}{1-\eps}$), no al límite. Con $\eps$ constante se ganan cifras a ritmo
geométrico: $\cO(\log(1/\delta))$ iteraciones, cada una con un \emph{sketch} de tamaño $m\asymp d$.
\textbf{Newton Sketch es exactamente esta idea aplicada a un hessiano que cambia con $x$.}
\end{ideaclave}

\begin{observacion}[Tasa exacta en el caso gaussiano]\label{obs:ihs-exacta}
Con \emph{sketches} gaussianos independientes en cada iteración, $W$ es de Wishart y se puede calcular
$\E\nrm{e_{t+1}}_2^2=\rho^2\nrm{e_t}_2^2$ con
\[ \rho^2=1-2\theta_1+\theta_2,\qquad\theta_1=\frac{m}{m-d-1},\quad\theta_2=\frac{m^2(m-1)}{(m-d)(m-d-1)(m-d-3)}, \]
usando $\E[W^{-1}]=\theta_1I$ y $\E[W^{-2}]=\theta_2I$. Para $m\gg d$, $\rho^2\approx d/m$: la tasa es
$\approx\sqrt{d/m}$ (un cálculo de Monte Carlo con $d=50$ confirma la fórmula con tres cifras). Con paso unitario
y $m=4d$ resulta $\rho\approx0.86$, bastante lento; por eso en la práctica se toma $m$ algo mayor o se añade un
tamaño de paso (Lacotte y Pilanci, 2020).
\end{observacion}

\section{Regresión \emph{ridge} y dimensión efectiva}

Para $\min_x\tfrac12\nrm{Ax-b}_2^2+\tfrac\mu2\nrm{x}_2^2$ el hessiano es $A\T A+\mu I$ y el IHS se convierte en
$x_{t+1}=x_t-(A\T S_t\T S_tA+\mu I)^{-1}\bigl(A\T(Ax_t-b)+\mu x_t\bigr)$. El mismo argumento del
teorema~\ref{teo:ihs}, ahora con la condición \eqref{eq:embed-reg}, da convergencia lineal; y por la
sección~\ref{sec:dim-efectiva} basta $m\asymp\dmu$ en lugar de $m\asymp d$. Lacotte y Pilanci (2020) propusieron
además un esquema \emph{adaptativo} que no necesita conocer $\dmu$: empieza con $m$ pequeño y lo duplica cuando la
convergencia observada es más lenta que la prevista. Es el antecedente directo del capítulo~\ref{cap:ans}.

\begin{figure}[ht]
\centering
\begin{tikzpicture}
\begin{semilogyaxis}[width=10.5cm,height=6.2cm,xlabel={iteración $t$},
  ylabel={$\nrm{A(x_t-x^\star)}_2/\nrm{Ax^\star}_2$},grid=both,grid style={gray!20},
  legend pos=south west,legend style={font=\footnotesize},tick label style={font=\footnotesize},label style={font=\small}]
\addplot[color=blue!70!black, mark=o, thick] coordinates {(0,1) (1,0.9199) (2,0.6019) (3,0.362) (4,0.2731) (5,0.1828) (6,0.1737) (7,0.1385) (8,0.116) (9,0.0995) (10,0.1092) (11,0.07914) (12,0.05888)};
\addlegendentry{$m=4d$}
\addplot[color=red!70!black, mark=square, thick] coordinates {(0,1) (1,0.2881) (2,0.1327) (3,0.04513) (4,0.02175) (5,0.009641) (6,0.004289) (7,0.002165) (8,0.0008951) (9,0.000397) (10,0.0001553) (11,5.724e-05) (12,2.536e-05)};
\addlegendentry{$m=8d$}
\addplot[color=green!50!black, mark=triangle, thick] coordinates {(0,1) (1,0.2428) (2,0.06341) (3,0.01902) (4,0.004835) (5,0.00154) (6,0.0005298) (7,0.0001405) (8,4.947e-05) (9,1.155e-05) (10,4.15e-06) (11,1.135e-06) (12,2.866e-07)};
\addlegendentry{$m=16d$}
\end{semilogyaxis}
\end{tikzpicture}
\caption{IHS con \emph{sketch} gaussiano renovado en cada iteración y paso unitario ($n=8192$, $d=50$, $\kappa(A)=10^6$).}
\label{fig:ihs}
\end{figure}

En la figura~\ref{fig:ihs} la convergencia es lineal y \emph{no depende de} $\kappa(A)=10^6$. Las tasas empíricas
(media geométrica entre las iteraciones 2 y 8) son $0.76$, $0.44$ y $0.30$ para $m=4d$, $8d$ y $16d$; la
observación~\ref{obs:ihs-exacta} predice $0.86$, $0.46$ y $0.29$ en media cuadrática.

\begin{resumen}
\begin{itemize}
  \item \emph{Sketch-and-solve}: una sola resolución pequeña; error relativo $\approx d/m$; solo baja precisión.
  \item \emph{Sketch-and-precondition}: $\kappa(AR^{-1})=\cO(1)$ con $m\asymp d$; alta precisión en $\cO(\log(1/\delta))$ iteraciones.
  \item IHS: \emph{sketch} del hessiano y gradiente exacto; convergencia lineal con factor $\eps/(1-\eps)$; es Newton Sketch para una cuadrática.
  \item Con regularización, el tamaño lo fija $\dmu$.
\end{itemize}
\end{resumen}

\chapter{Aplicación a sistemas lineales y a programas lineales}\label{cap:lin}

«Problemas lineales» puede significar dos cosas, y el \emph{sketching} interviene en ambas: \emph{sistemas de
ecuaciones lineales} $Ax=b$, donde el \emph{sketch} da lugar a métodos iterativos aleatorizados
(\emph{sketch-and-project}), y \emph{programas lineales}, donde el \emph{sketch} acelera el paso de Newton del
método de barrera.

\section{Sistemas lineales: \emph{sketch-and-project}}

Sea $Ax=b$ un sistema \emph{consistente} con $A\in\R^{n\times d}$ de rango $d$ y solución $x^\star$. La idea de
Gower y Richtárik (2015) es no resolver el sistema entero, sino una versión comprimida distinta en cada
iteración, y proyectar sobre ella.

\begin{definicion}[Iteración \emph{sketch-and-project}]\label{def:sap}
Sea $\mathsf B\succ0$ una matriz $d\times d$ (define la norma $\nrm{x}_{\mathsf B}^2=x\T\mathsf Bx$) y $\mathcal D$ una
distribución sobre matrices $S\in\R^{n\times q}$. En cada iteración se extrae $S\sim\mathcal D$ y
\begin{equation}\label{eq:sap}
  x_{k+1}=\argmin_{x}\;\nrm{x-x_k}_{\mathsf B}^2\quad\text{s.a.}\quad S\T Ax=S\T b .
\end{equation}
\end{definicion}

Aquí el \emph{sketch} actúa sobre las \emph{ecuaciones}: $S\T A\in\R^{q\times d}$ son $q$ combinaciones lineales de
las $n$ ecuaciones, con $q$ muy pequeño (a menudo $q=1$). La solución de \eqref{eq:sap} es explícita:
\begin{equation}\label{eq:sap-explicita}
  x_{k+1}=x_k-\mathsf B^{-1}A\T S\,(S\T A\mathsf B^{-1}A\T S)^{\dagger}\,S\T(Ax_k-b).
\end{equation}

\begin{teorema}[Convergencia lineal en media {\cite{GR15}}]\label{teo:sap}
Sea $Z=A\T S(S\T A\mathsf B^{-1}A\T S)^{\dagger}S\T A$ y supóngase $\E[Z]\succ0$. Entonces
\[ \E\nrm{x_k-x^\star}_{\mathsf B}^2\;\le\;\rho^{\,k}\,\nrm{x_0-x^\star}_{\mathsf B}^2,\qquad
   \rho=1-\lambda_{\min}\bigl(\mathsf B^{-1/2}\,\E[Z]\,\mathsf B^{-1/2}\bigr)<1 . \]
\end{teorema}

\begin{proof}
Como $Ax^\star=b$, \eqref{eq:sap-explicita} da $x_{k+1}-x^\star=(I-\mathsf B^{-1}Z)(x_k-x^\star)$. Sea
$e_k=\mathsf B^{1/2}(x_k-x^\star)$ y $\Pi=\mathsf B^{-1/2}Z\mathsf B^{-1/2}$; entonces $e_{k+1}=(I-\Pi)e_k$. La matriz $\Pi$ es
el proyector ortogonal sobre la imagen de $\mathsf B^{-1/2}A\T S$ (es de la forma $C(C\T C)^{\dagger}C\T$), así que
$\nrm{e_{k+1}}_2^2=\nrm{e_k}_2^2-e_k\T\Pi e_k$. Tomando esperanza condicionada a $x_k$,
$\E[e_k\T\Pi e_k]=e_k\T\E[\Pi]e_k\ge\lambda_{\min}(\E[\Pi])\nrm{e_k}_2^2$.
\end{proof}

La estructura es transparente: cada iteración \emph{proyecta ortogonalmente} el error sobre un subespacio
aleatorio; ninguna proyección aislada lo anula, pero \emph{en media} todas las direcciones se contraen.

\subsection{Casos particulares clásicos}

\begin{ejemplo}[Kaczmarz aleatorizado]\label{ej:kaczmarz}
Con $\mathsf B=I$ y $S=e_i$ (una ecuación al azar), \eqref{eq:sap-explicita} es
\[ x_{k+1}=x_k-\frac{a_i\T x_k-b_i}{\nrm{a_i}_2^2}\,a_i : \]
la proyección de $x_k$ sobre el hiperplano de la ecuación $i$. Si $i$ se elige con probabilidad
$p_i=\nrm{a_i}_2^2/\nrm{A}_F^2$, entonces $\E[Z]=\sum_ip_i\,a_ia_i\T/\nrm{a_i}_2^2=A\T A/\nrm{A}_F^2$ y
\[ \rho=1-\frac{\sigma_{\min}(A)^2}{\nrm{A}_F^2}, \]
que es el teorema de Strohmer y Vershynin (2009). Cada iteración cuesta $\cO(d)$.
\end{ejemplo}

\begin{ejemplo}[Descenso por coordenadas]\label{ej:cd}
Si $A$ es simétrica definida positiva ($n=d$), $\mathsf B=A$ y $S=e_i$ con $p_i=A_{ii}/\tr(A)$, se obtiene el descenso
por coordenadas aleatorizado, con $\rho=1-\lambda_{\min}(A)/\tr(A)$ (Leventhal y Lewis, 2010).
\end{ejemplo}

\begin{ejemplo}[Versiones por bloques y gaussianas]
Con $S$ formada por $q$ columnas de la identidad se obtienen Kaczmarz y descenso por coordenadas por bloques;
con $S$ gaussiana, variantes con direcciones aleatorias densas. Aumentar $q$ mejora $\rho$ a cambio de resolver
un sistema $q\times q$ en cada iteración.
\end{ejemplo}

\begin{cuidado}[Dos usos distintos del \emph{sketch} para sistemas lineales]
\begin{itemize}
  \item \textbf{\emph{Sketch-and-project}}: muchos \emph{sketches} diminutos, uno por iteración. La tasa $\rho$
        \emph{depende del condicionamiento} (en Kaczmarz, de $\nrm{A}_F^2/\sigma_{\min}^2$). Útil cuando solo se puede acceder a unas pocas filas cada vez.
  \item \textbf{\emph{Sketch-and-precondition}} (capítulo~\ref{cap:ls}): un \emph{sketch} grande ($m\asymp d$) que
        elimina la dependencia del condicionamiento. Es el que usa Newton Sketch.
\end{itemize}
\end{cuidado}

\subsection{El sistema de Newton como problema de mínimos cuadrados}

El sistema lineal que más nos interesa es el de Newton, $H v=-g$ con $H=B\T B$. No es un sistema cualquiera:
es la ecuación normal del problema
\[ \min_{v}\;\tfrac12\nrm{Bv}_2^2+\ip{g}{v} . \]
Hay entonces dos formas de explotar el \emph{sketch} $SB$, que corresponden a las estrategias segunda y
tercera del capítulo~\ref{cap:ls}:
\begin{enumerate}
  \item \textbf{Newton Sketch}: reemplazar $H$ por $H_S=(SB)\T(SB)$ y resolver $H_Sv=-g$ exactamente. El paso es
        inexacto, pero el gradiente es exacto (como en el IHS).
  \item \textbf{Newton con gradiente conjugado precondicionado}: resolver $Hv=-g$ con gradiente conjugado usando $H_S$
        como precondicionador. Si $(1-\eps)H\preceq H_S\preceq(1+\eps)H$, entonces
        $\kappa(H_S^{-1/2}HH_S^{-1/2})\le\frac{1+\eps}{1-\eps}$ y bastan unas pocas iteraciones, cada una con un producto
        por $B$ y otro por $B\T$, de coste $\cO(nd)$.
\end{enumerate}
Ambas evitan formar $B\T B$.

\section{Programas lineales}\label{sec:pl}

Considérese el programa lineal en forma de desigualdades
\begin{equation}\label{eq:pl}
  \min_{x\in\R^d}\;c\T x\quad\text{s.a.}\quad Ax\le b,\qquad A\in\R^{n\times d},
\end{equation}
con muchas más restricciones que variables ($n\gg d$). El método de barrera (algoritmo~\ref{alg:barrera})
minimiza, para $\theta$ creciente,
\[ F_\theta(x)=\theta\,c\T x-\sum_{i=1}^n\log(b_i-a_i\T x), \]
que es autoconcordante estándar. Con holguras $s_i=b_i-a_i\T x>0$:
\[ \nabla F_\theta(x)=\theta c+A\T(1/s),\qquad\nabla^2F_\theta(x)=A\T\diag(1/s_i^2)A=B\T B,\qquad B=\diag(1/s_i)\,A . \]
Cada paso de Newton cuesta $\cO(nd^2)$. Con Newton Sketch se forma $SB$ ---con SRHT, en $\cO(nd\log m)$--- y
el paso cuesta $\cO(nd\log m+md^2)$; con $m\asymp d$, $\cO(nd\log d+d^3)$. Lo mismo vale para programas
cuadráticos convexos, añadiendo $\theta Q$ al hessiano sin aplicarle \emph{sketch}.

\begin{cuidado}[La barrera se vuelve coherente cerca del óptimo]
En un vértice óptimo no degenerado hay $d$ restricciones activas; a lo largo de la trayectoria central sus
holguras tienden a $0$ y las filas correspondientes de $B=\diag(1/s)A$ crecen sin límite: \emph{$d$ filas
concentran casi toda la palanca}. Es exactamente la «matriz coherente» del experimento de la
sección~\ref{sec:exp-embed}. Consecuencias:
\begin{itemize}
  \item el muestreo uniforme de restricciones fracasa al final del método de barrera;
  \item los \emph{sketches} ajenos a los datos que mezclan (gaussiano, SRHT) siguen funcionando con $m\asymp d$;
  \item CountSketch necesita $m\asymp d^2$ para no mezclar dos restricciones activas en la misma cubeta.
\end{itemize}
\end{cuidado}

\begin{resumen}
\begin{itemize}
  \item \emph{Sketch-and-project}: proyectar sobre una versión comprimida aleatoria del sistema; convergencia
        lineal en media con $\rho=1-\lambda_{\min}(\mathsf B^{-1/2}\E[Z]\mathsf B^{-1/2})$. Kaczmarz y descenso por coordenadas son casos particulares.
  \item El sistema de Newton es un problema de mínimos cuadrados con matriz $B$; se le aplica todo el capítulo~\ref{cap:ls}.
  \item En programación lineal con barrera, $B=\diag(1/s)A$; Newton Sketch reduce el coste por paso de $nd^2$ a $nd\log d+d^3$.
  \item Cerca del óptimo $B$ es coherente: hacen falta \emph{sketches} que mezclen.
\end{itemize}
\end{resumen}

\part{Newton Sketch}
\chapter{Newton Sketch}\label{cap:ns}

Este capítulo presenta el algoritmo de Pilanci y Wainwright~\cite{PW17}: el método de Newton en el que el
hessiano se sustituye por una versión comprimida aleatoriamente. Es la confluencia de los capítulos
anteriores: el encaje de subespacio (capítulo~\ref{cap:embed}) controla el error del hessiano y la
autoconcordancia (capítulo~\ref{cap:sc}) convierte ese control en una garantía de convergencia sin constantes
del problema.

\section{Derivación}

Sea $f$ convexa y dos veces diferenciable con raíz del hessiano $B(x)=\nabla^2 f(x)^{1/2}\in\R^{n\times d}$
(sección~\ref{sec:raiz-hess}), y $\cC\subseteq\R^d$ un conjunto convexo cerrado. El paso de Newton (proyectado) es
un problema de mínimos cuadrados con término lineal:
\[ x_{t+1}^{\mathrm{Newton}}=\argmin_{x\in\cC}\Bigl\{\tfrac12\nrm{B(x_t)(x-x_t)}_2^2+\ip{\nabla f(x_t)}{x-x_t}\Bigr\}. \]

\begin{definicion}[Paso de Newton Sketch]\label{def:ns}
Dada una matriz de \emph{sketch} $S_t\in\R^{m\times n}$, \emph{independiente en cada iteración},
\begin{equation}\label{eq:ns}
  x_{t+1}=\argmin_{x\in\cC}\Bigl\{\tfrac12\nrm{S_tB(x_t)(x-x_t)}_2^2+\ip{\nabla f(x_t)}{x-x_t}\Bigr\}.
\end{equation}
En el caso sin restricciones ($\cC=\R^d$):
\begin{equation}\label{eq:ns-unc}
  x_{t+1}=x_t+s\,\vnsk,\qquad\vnsk=-H_S(x_t)^{-1}\nabla f(x_t),\qquad H_S(x)=B(x)\T S\T SB(x),
\end{equation}
con paso $s=1$ o elegido por búsqueda lineal.
\end{definicion}

Tres observaciones de diseño:
\begin{enumerate}
  \item \textbf{Solo se comprime el hessiano; el gradiente es exacto.} Como en el IHS (sección~\ref{sec:ihs}), el
        punto fijo es $\nabla f(x)=0$: el método converge a la solución \emph{exacta}. Comprimir los datos una vez y
        optimizar sobre los datos comprimidos daría un método sesgado.
  \item \textbf{No se forma $B\T B$.} Se forma $S_tB\in\R^{m\times d}$ y luego $(S_tB)\T(S_tB)$, de coste $\cO(md^2)$.
  \item \textbf{\emph{Sketch} nuevo en cada iteración.} Así $S_t$ es independiente de $x_t$ y se puede aplicar la
        garantía de encaje a la matriz fija $B(x_t)$; además los errores no se acumulan de forma coherente.
\end{enumerate}

\begin{definicion}[Decremento de Newton aproximado]
\[ \lamt(x)=\bigl(\nabla f(x)\T H_S(x)^{-1}\nabla f(x)\bigr)^{1/2}=\bigl(-\ip{\nabla f(x)}{\vnsk}\bigr)^{1/2}. \]
\end{definicion}

Es calculable sin coste adicional (es un producto escalar con el paso ya calculado) y sustituye a $\lamf$
tanto en la condición de Armijo como en el criterio de parada.

\begin{algorithm}[ht]
\caption{Newton Sketch con búsqueda lineal (sin restricciones)}\label{alg:ns}
\begin{algorithmic}[1]
\Require $x_0\in\dom f$, tamaño $m$, parámetros $(a,b)$, tolerancia $\delta$
\For{$t=0,1,2,\dots$}
  \State Generar $S_t\in\R^{m\times n}$ independiente; formar $S_tB(x_t)$ y $H_S=(S_tB)\T(S_tB)$
  \State $\vnsk\gets-H_S^{-1}\nabla f(x_t)$;\quad$\lamt^2\gets-\ip{\nabla f(x_t)}{\vnsk}$
  \State \textbf{si} $\lamt^2/2\le\delta$ \textbf{entonces devolver} $x_t$
  \State $s\gets1$;\quad\textbf{mientras} $f(x_t+s\vnsk)>f(x_t)-a\,s\,\lamt^2$ \textbf{hacer} $s\gets b\,s$
  \State $x_{t+1}\gets x_t+s\,\vnsk$
\EndFor
\end{algorithmic}
\end{algorithm}

\section{El lema básico: direcciones y decrementos cercanos}

\begin{lema}[Cercanía del paso con \emph{sketch}]\label{lem:ns-cerca}
Sea $H=B\T B\succ0$ y supóngase que $S$ es un encaje de subespacio con distorsión $\eps\in(0,1)$ para $B$, es
decir, $(1-\eps)H\preceq H_S\preceq(1+\eps)H$. Entonces, con $\vne=-H^{-1}g$ y $\vnsk=-H_S^{-1}g$:
\begin{align}
  \nrm{\vnsk-\vne}_x\;&\le\;\frac{\eps}{1-\eps}\,\nrm{\vne}_x=\frac{\eps}{1-\eps}\,\lamf(x), \label{eq:ns-dir}\\
  \frac{\lamf(x)^2}{1+\eps}\;&\le\;\lamt(x)^2\;\le\;\frac{\lamf(x)^2}{1-\eps}, \label{eq:ns-dec}\\
  \nrm{\vnsk}_x^2\;&\le\;\frac{\lamt(x)^2}{1-\eps}. \label{eq:ns-norma}
\end{align}
\end{lema}

\begin{proof}
Sea $B=UT$ con $U\in\R^{n\times d}$ de columnas ortonormales y $T$ invertible, $W=U\T S\T SU$ y $z=T^{-\top}g$.
Entonces $H=T\T T$, $H_S=T\T WT$, y
\[ B\vne=-UT(T\T T)^{-1}g=-Uz,\qquad B\vnsk=-UT(T\T WT)^{-1}g=-UW^{-1}z . \]
Por tanto $\nrm{\vne}_x=\nrm{z}_2=\lamf(x)$, $\nrm{\vnsk-\vne}_x=\nrm{(W^{-1}-I)z}_2$, $\lamt(x)^2=z\T W^{-1}z$ y
$\nrm{\vnsk}_x^2=z\T W^{-2}z$. Como los valores propios de $W$ están en $[1-\eps,1+\eps]$
(proposición~\ref{prop:embed-equiv}), los de $W^{-1}-I$ están en $[-\frac{\eps}{1+\eps},\frac{\eps}{1-\eps}]$, los de
$W^{-1}$ en $[\frac{1}{1+\eps},\frac{1}{1-\eps}]$, y $W^{-2}\preceq\frac{1}{1-\eps}W^{-1}$.
\end{proof}

\begin{ideaclave}
En las coordenadas adecuadas, el paso exacto es $z$ y el paso con \emph{sketch} es $W^{-1}z$, con
$W=U\T S\T SU\approx I$. \textbf{El error del paso es relativo}: proporcional al tamaño del propio paso, no a
$\kappa$ ni a $n$. Un error relativo constante en la dirección de Newton degrada la convergencia cuadrática a
lineal, pero conserva la propiedad esencial: una tasa que no depende del condicionamiento.
\end{ideaclave}

\section{Convergencia local para funciones fuertemente convexas y suaves}

El primer resultado de~\cite{PW17} usa las hipótesis clásicas del capítulo~\ref{cap:opt}. Sean $\mu$ y $L$ los
valores propios mínimo y máximo de $\nabla^2 f(x^\star)$ (restringidos al cono tangente $\cK$ de $\cC$ en $x^\star$
si hay restricciones) y $\LH$ la constante de Lipschitz del hessiano.

\begin{teorema}[Convergencia local lineal-cuadrática {\cite[Teorema~1]{PW17}}]\label{teo:pw1}
Sea $\eps\in(0,1)$ y $\nrm{x_0-x^\star}_2\le\dfrac{\mu}{8\LH}$. Si el tamaño del \emph{sketch} (subgaussiano) cumple
\[ m\;\ge\;\frac{c}{\eps^2}\,\max_{x\in\cC}\cW^2\bigl(B(x)\cK\bigr), \]
entonces, con probabilidad al menos $1-c_1Ne^{-c_2m}$, las $N$ primeras iteraciones de \eqref{eq:ns} con paso
unitario satisfacen
\begin{equation}\label{eq:lin-cuad}
  \nrm{x_{t+1}-x^\star}_2\;\le\;\eps\,\frac{L}{\mu}\,\nrm{x_t-x^\star}_2+\frac{4\LH}{\mu}\,\nrm{x_t-x^\star}_2^2 .
\end{equation}
Aquí $c,c_1,c_2$ son constantes universales y $\cW$ es el ancho gaussiano (véase abajo).
\end{teorema}

\paragraph{Lectura.} La recurrencia \eqref{eq:lin-cuad} es \emph{lineal-cuadrática}. El término cuadrático es
el de Newton exacto. El término lineal, $\eps\kappa\nrm{x_t-x^\star}$, es el precio del \emph{sketch}. Lejos del
óptimo domina el cuadrático y el método se comporta como Newton; cerca domina el lineal y la convergencia es
geométrica con razón $\eps\kappa$. Bajo la condición inicial, $\frac{4\LH}{\mu}\nrm{x_t-x^\star}\le\frac12$, de modo
que $\nrm{x_N-x^\star}\le\frac{\mu}{8\LH}(\frac12+\eps\kappa)^N$. Haciendo $\eps_t\to0$ (aumentando $m$ con $t$) se
obtiene convergencia superlineal.

\begin{cuidado}[Por qué este teorema no es la última palabra]
La contracción exige $\eps\kappa<\tfrac12$, es decir $m\gtrsim\kappa^2d$: ¡reaparece el número de condición! No es
un defecto del método sino del análisis en norma euclídea (el mismo defecto del
teorema~\ref{teo:newton-clasico}). El lema~\ref{lem:ns-cerca} dice que en la norma local el error relativo es
$\eps/(1-\eps)$ \emph{sin} $\kappa$. El análisis autoconcordante de la sección siguiente trabaja en esa norma.
\end{cuidado}

\paragraph{Ancho gaussiano y restricciones.} El ancho gaussiano de un conjunto $\mathcal L$ es
\[ \cW(\mathcal L)=\E_g\,\sup_{z\in\mathcal L}\,\abs{\ip{g}{z}},\qquad g\sim N(0,I), \]
y mide el «tamaño» de $\mathcal L$ visto por una proyección aleatoria. En el teorema se aplica a $B(x)\cK\cap\mathbb S^{n-1}$. Sin restricciones, $B\cK$ es un subespacio
de dimensión $d$ y $\cW^2\le d$: se recupera $m\gtrsim d/\eps^2$. Con restricciones que inducen estructura (por
ejemplo una bola $\ell_1$ con solución $k$-dispersa), el cono tangente es mucho más pequeño y
$\cW^2\asymp k\log d$: basta un \emph{sketch} de tamaño proporcional a la \emph{dimensión intrínseca} del problema.
Para la SRHT valen cotas análogas con factores logarítmicos adicionales.

\section{Convergencia global para funciones autoconcordantes}

Sea ahora $f$ autoconcordante estándar y estrictamente convexa, sin restricciones. Se fija $\eps$ y los
parámetros de la búsqueda lineal de modo que, con $\varrho=\frac{1+\eps}{1-\eps}$, se cumpla
$a<1-\tfrac12\varrho^2$, y se definen
\begin{equation}\label{eq:eta-nu}
  \eta=\frac18\cdot\frac{1-\frac12\varrho^2-a}{\varrho^{3}},\qquad\nu=a\,b\,\frac{\eta^2}{1+\varrho\,\eta}.
\end{equation}

\begin{lema}[Las dos fases {\cite[Lema~3]{PW17}}]\label{lem:pw-fases}
Supóngase que en la iteración $t$ el \emph{sketch} satisface la cercanía \eqref{eq:ns-dir} con parámetro $\eps$.
\begin{enumerate}[label=(\alph*)]
  \item \textbf{Fase 1.} Si $\lamt(x_t)>\eta$, entonces $f(x_{t+1})-f(x_t)\le-\nu$.
  \item \textbf{Fase 2.} Si $\lamt(x_t)\le\eta$, la búsqueda lineal acepta $s=1$ y
        $\lamf(x_{t+1})\le\tfrac{16}{25}\,\lamf(x_t)$.
\end{enumerate}
\end{lema}

\begin{proof}[Esquema]
(a) Sea $\lambda=\lamt(x_t)$ y $c$ tal que $\nrm{\vnsk}_x\le c\lambda$ (por \eqref{eq:ns-norma} sirve
$c=(1-\eps)^{-1/2}$; en~\cite{PW17} se usa $c=\varrho$). Por el teorema~\ref{teo:sc-cotas}(iv),
$f(x+s\vnsk)\le f(x)-s\lambda^2+\omega_*(sc\lambda)$ mientras $sc\lambda<1$. Igual que en el caso exacto
(teorema~\ref{teo:sc-damped}), para $\hat s=1/(1+c\lambda)$ el lado derecho queda por debajo de $f(x)-a\hat s\lambda^2$, de
modo que el retroceso termina con $s\ge b/(1+c\lambda)$ (aquí se usa $\log(1+u)\le u-\frac{u^2}{2(1+u)}$ y la condición $a\le1-\frac12c^2$) y el descenso es al menos $ab\lambda^2/(1+c\lambda)$, que es creciente
en $\lambda$ y por tanto $\ge\nu$ cuando $\lambda>\eta$.
(b) Con paso unitario, se repite la demostración del teorema~\ref{teo:sc-quad} con $\vnsk=\vne+(\vnsk-\vne)$: el
primer sumando aporta el término cuadrático $\cO(\lamf^2)$ y el segundo, por \eqref{eq:ns-dir}, un término lineal
$\cO(\eps\lamf)$. Para $\lamf$ por debajo del umbral, la suma es $\le\frac{16}{25}\lamf$.
\end{proof}

\begin{teorema}[Complejidad de Newton Sketch {\cite[Teorema~2]{PW17}}]\label{teo:pw2}
Sea $f$ autoconcordante estándar y estrictamente convexa. Si el \emph{sketch} (subgaussiano o SRHT) tiene tamaño
$m\ge c_3\,d/\eps^2$ (con factores logarítmicos adicionales para la SRHT), el algoritmo~\ref{alg:ns} devuelve un punto con
$f(x)-f^\star\le\delta$ en a lo sumo
\[ N=\frac{f(x_0)-f^\star}{\nu}+0.65\,\log_2\frac{1}{16\,\delta} \]
iteraciones, con probabilidad al menos $1-c_1Ne^{-c_2m}$.
\end{teorema}

\paragraph{Lectura.}
\begin{itemize}
  \item La estructura es la del teorema~\ref{teo:sc-complejidad}: fase de descenso constante más fase rápida. El
        cambio es que $\log_2\log_2(1/\delta)$ pasa a ser un término proporcional a $\log_2(1/\delta)$:
        \textbf{convergencia lineal}, con razón $\frac{16}{25}$ en $\lamf$.
  \item \textbf{No aparece $\kappa$}, ni $\mu$, ni $L$: solo $a$, $b$ y $\eps$. Y $m$ depende solo de $d$.
  \item La probabilidad de fallo se controla con la cota de la unión sobre las $N$ iteraciones; como decae
        exponencialmente en $m$, no es una restricción en la práctica.
  \item Las constantes \eqref{eq:eta-nu} son muy conservadoras. Con $\eps=\tfrac18$, $a=0.1$, $b=0.5$ se obtiene
        $\varrho=\tfrac97$, $\eta\approx4.3\cdot10^{-3}$ y $\nu\approx9.3\cdot10^{-7}$. En los experimentos la fase 1 dura
        unas pocas iteraciones, no $10^6(f(x_0)-f^\star)$.
\end{itemize}

\section{Coste y comparación}

Con SRHT y $m\asymp d$, una iteración cuesta
\[ \underbrace{\cO(nd)}_{B\text{ y }\nabla f}+\underbrace{\cO(nd\log m)}_{S_tB}+\underbrace{\cO(md^2)}_{(S_tB)\T(S_tB)}+\underbrace{\cO(d^3)}_{\text{sistema}}=\cO(nd\log d+d^3), \]
que es $\cO(nd\log d)$, casi lineal en el tamaño de los datos, cuando $n\gtrsim d^2$.

\begin{table}[ht]
\centering\small
\caption{Coste total para alcanzar precisión $\delta$ en un problema del tipo \eqref{eq:problema-general} con
$n\gtrsim d^2$ (órdenes de magnitud, según~\cite{PW17,LWP21}).}
\label{tab:complejidad-ns}
\begin{tabular}{@{}lll@{}}
\toprule
Método & Coste total & ¿Depende de $\kappa$? \\
\midrule
Descenso de gradiente & $\kappa\,nd\,\log(1/\delta)$ & sí \\
Gradiente acelerado (Nesterov) & $\sqrt\kappa\,nd\,\log(1/\delta)$ & sí \\
Newton exacto & $nd^2\,\log\log(1/\delta)$ & no \\
Newton Sketch, $m\asymp d$ & $nd\log d\,\log(1/\delta)$ & no \\
\bottomrule
\end{tabular}
\end{table}

Newton Sketch cambia el factor $d$ de Newton por $\log d\cdot\log(1/\delta)/\log\log(1/\delta)$: gana cuando $d$ es grande. El
precio es pasar de convergencia cuadrática a lineal. Recuperar la cuadrática es uno de los objetivos del
capítulo~\ref{cap:ans}.

\section{Restricciones, \emph{sketch} parcial y método de barrera}

\paragraph{\emph{Sketch} parcial.} Si $f=f_0+g$ y el hessiano de $g$ es barato (diagonal, por ejemplo), solo se
comprime $f_0$:
\begin{equation}\label{eq:ns-parcial}
  H_S(x)=B_0(x)\T S\T SB_0(x)+\nabla^2g(x),\qquad B_0=\nabla^2f_0^{1/2}.
\end{equation}
El lema~\ref{lem:ns-cerca} sigue valiendo, porque $(1-\eps)B_0\T B_0\preceq B_0\T S\T SB_0\preceq(1+\eps)B_0\T B_0$ implica
$(1-\eps)H\preceq H_S\preceq(1+\eps)H$ al sumar $\nabla^2g\succeq0$ a los tres miembros. Esta es la forma que explota
el capítulo siguiente, donde además se aprovecha que $\nabla^2 g\succeq\mu I$ para \emph{reducir} $m$.

\paragraph{Método de barrera.} Para problemas con restricciones, \cite{PW17} usa Newton Sketch como
resolvedor del centrado en el algoritmo~\ref{alg:barrera}: cada $F_\theta$ es autoconcordante y se le aplica el
teorema~\ref{teo:pw2}. El número total de pasos de Newton Sketch es el número de centrados,
$\lceil\log(n_c/(\theta_0\delta))/\log\mu_b\rceil$ con $n_c$ restricciones, multiplicado por una cota por centrado de la
forma $\frac{n_c(\mu_b-1-\log\mu_b)}{\nu}+0.65\log_2\frac{1}{16\delta}$ \cite[Teorema~4]{PW17}. Para un programa lineal
con $n$ restricciones y $d$ variables, cada paso cuesta $\cO(nd\log d+d^3)$ en lugar de $\cO(nd^2)$
(sección~\ref{sec:pl}).

\begin{resumen}
\begin{itemize}
  \item Newton Sketch: $\vnsk=-(B\T S_t\T S_tB)^{-1}\nabla f$; gradiente exacto, hessiano comprimido, $S_t$ nuevo en cada iteración.
  \item Bajo un encaje con distorsión $\eps$: $\nrm{\vnsk-\vne}_x\le\frac{\eps}{1-\eps}\lamf$ y $\lamt^2\in[\frac{1}{1+\eps},\frac{1}{1-\eps}]\lamf^2$.
  \item Análisis euclídeo: recurrencia lineal-cuadrática $\eps\kappa\nrm{\Delta}+\frac{4\LH}{\mu}\nrm{\Delta}^2$ (local, con $\kappa$).
  \item Análisis autoconcordante: $(f(x_0)-f^\star)/\nu+0.65\log_2(1/(16\delta))$ iteraciones con $m\asymp d$; sin $\kappa$.
  \item Coste total $\cO(nd\log d\log(1/\delta))$: lineal en los datos salvo logaritmos.
\end{itemize}
\end{resumen}


\begin{thebibliography}{99}
\addcontentsline{toc}{chapter}{Bibliografía}
\small

\bibitem{PW17} M. Pilanci y M. J. Wainwright. Newton Sketch: A near linear-time optimization algorithm with
linear-quadratic convergence. \emph{SIAM Journal on Optimization}, 27(1):205--245, 2017. (Los enunciados de este
documento siguen la versión de arXiv:1505.02250.)

\bibitem{LWP21} J. Lacotte, Y. Wang y M. Pilanci. Adaptive Newton Sketch: Linear-time optimization with quadratic
convergence and effective Hessian dimensionality. \emph{Proceedings of the 38th International Conference on Machine
Learning (ICML)}, PMLR 139, 2021. arXiv:2105.07291.

\bibitem{PW16} M. Pilanci y M. J. Wainwright. Iterative Hessian sketch: Fast and accurate solution approximation for
constrained least-squares. \emph{Journal of Machine Learning Research}, 17(53):1--38, 2016.

\bibitem{LP20} J. Lacotte y M. Pilanci. Effective dimension adaptive sketching methods for faster regularized
least-squares optimization. \emph{Advances in Neural Information Processing Systems (NeurIPS)}, 2020.

\bibitem{BV} S. Boyd y L. Vandenberghe. \emph{Convex Optimization}. Cambridge University Press, 2004.

\bibitem{Nesterov} Y. Nesterov. \emph{Lectures on Convex Optimization}, 2.ª ed. Springer, 2018.

\bibitem{NN94} Y. Nesterov y A. Nemirovskii. \emph{Interior-Point Polynomial Algorithms in Convex Programming}. SIAM, 1994.

\bibitem{Bach10} F. Bach. Self-concordant analysis for logistic regression. \emph{Electronic Journal of Statistics},
4:384--414, 2010.

\bibitem{ZX15} Y. Zhang y L. Xiao. DiSCO: Distributed optimization for self-concordant empirical loss.
\emph{Proceedings of the 32nd International Conference on Machine Learning (ICML)}, 2015.

\bibitem{Woodruff14} D. P. Woodruff. Sketching as a tool for numerical linear algebra. \emph{Foundations and Trends in
Theoretical Computer Science}, 10(1--2):1--157, 2014.

\bibitem{MT20} P.-G. Martinsson y J. A. Tropp. Randomized numerical linear algebra: Foundations and algorithms.
\emph{Acta Numerica}, 29:403--572, 2020.

\bibitem{Mahoney11} M. W. Mahoney. Randomized algorithms for matrices and data. \emph{Foundations and Trends in Machine
Learning}, 3(2):123--224, 2011.

\bibitem{JL84} W. B. Johnson y J. Lindenstrauss. Extensions of Lipschitz mappings into a Hilbert space.
\emph{Contemporary Mathematics}, 26:189--206, 1984.

\bibitem{Sarlos06} T. Sarlós. Improved approximation algorithms for large matrices via random projections.
\emph{47th IEEE Symposium on Foundations of Computer Science (FOCS)}, 2006.

\bibitem{AC09} N. Ailon y B. Chazelle. The fast Johnson--Lindenstrauss transform and approximate nearest neighbors.
\emph{SIAM Journal on Computing}, 39(1):302--322, 2009.

\bibitem{Tropp11} J. A. Tropp. Improved analysis of the subsampled randomized Hadamard transform. \emph{Advances in
Adaptive Data Analysis}, 3(1--2):115--126, 2011.

\bibitem{BG13} C. Boutsidis y A. Gittens. Improved matrix algorithms via the subsampled randomized Hadamard transform.
\emph{SIAM Journal on Matrix Analysis and Applications}, 34(3):1301--1340, 2013.

\bibitem{Tropp12} J. A. Tropp. User-friendly tail bounds for sums of random matrices. \emph{Foundations of Computational
Mathematics}, 12:389--434, 2012.

\bibitem{CW13} K. L. Clarkson y D. P. Woodruff. Low rank approximation and regression in input sparsity time.
\emph{45th ACM Symposium on Theory of Computing (STOC)}, 2013.

\bibitem{NN13} J. Nelson y H. L. Nguyên. OSNAP: Faster numerical linear algebra algorithms via sparser subspace
embeddings. \emph{54th IEEE Symposium on Foundations of Computer Science (FOCS)}, 2013.

\bibitem{MM13} X. Meng y M. W. Mahoney. Low-distortion subspace embeddings in input-sparsity time and applications to
robust linear regression. \emph{45th ACM Symposium on Theory of Computing (STOC)}, 2013.

\bibitem{NN14} J. Nelson y H. L. Nguyên. Lower bounds for oblivious subspace embeddings. \emph{International
Colloquium on Automata, Languages and Programming (ICALP)}, 2014.

\bibitem{Cohen16} M. B. Cohen. Nearly tight oblivious subspace embeddings by trace inequalities. \emph{27th ACM-SIAM
Symposium on Discrete Algorithms (SODA)}, 2016.

\bibitem{CNW16} M. B. Cohen, J. Nelson y D. P. Woodruff. Optimal approximate matrix product in terms of stable rank.
\emph{International Colloquium on Automata, Languages and Programming (ICALP)}, 2016.

\bibitem{ACW17} H. Avron, K. L. Clarkson y D. P. Woodruff. Sharper bounds for regularized data fitting.
\emph{APPROX/RANDOM}, 2017.

\bibitem{DMMS11} P. Drineas, M. W. Mahoney, S. Muthukrishnan y T. Sarlós. Faster least squares approximation.
\emph{Numerische Mathematik}, 117(2):219--249, 2011.

\bibitem{RT08} V. Rokhlin y M. Tygert. A fast randomized algorithm for overdetermined linear least-squares regression.
\emph{Proceedings of the National Academy of Sciences}, 105(36):13212--13217, 2008.

\bibitem{AMT10} H. Avron, P. Maymounkov y S. Toledo. Blendenpik: Supercharging LAPACK's least-squares solver.
\emph{SIAM Journal on Scientific Computing}, 32(3):1217--1236, 2010.

\bibitem{MSM14} X. Meng, M. A. Saunders y M. W. Mahoney. LSRN: A parallel iterative solver for strongly over- or
underdetermined systems. \emph{SIAM Journal on Scientific Computing}, 36(2):C95--C118, 2014.

\bibitem{SV09} T. Strohmer y R. Vershynin. A randomized Kaczmarz algorithm with exponential convergence.
\emph{Journal of Fourier Analysis and Applications}, 15:262--278, 2009.

\bibitem{LL10} D. Leventhal y A. S. Lewis. Randomized methods for linear constraints: Convergence rates and
conditioning. \emph{Mathematics of Operations Research}, 35(3):641--654, 2010.

\bibitem{GR15} R. M. Gower y P. Richtárik. Randomized iterative methods for linear systems. \emph{SIAM Journal on
Matrix Analysis and Applications}, 36(4):1660--1690, 2015.

\bibitem{RKM19} F. Roosta-Khorasani y M. W. Mahoney. Sub-sampled Newton methods. \emph{Mathematical Programming},
174:293--326, 2019.

\bibitem{EM15} M. A. Erdogdu y A. Montanari. Convergence rates of sub-sampled Newton methods. \emph{Advances in Neural
Information Processing Systems (NeurIPS)}, 2015.

\bibitem{WW} A. J. Wood, B. F. Wollenberg y G. B. Sheblé. \emph{Power Generation, Operation, and Control}, 3.ª ed.
Wiley, 2014.

\bibitem{PP91} M. V. F. Pereira y L. M. V. G. Pinto. Multi-stage stochastic optimization applied to energy planning.
\emph{Mathematical Programming}, 52:359--375, 1991.

\bibitem{SDR} A. Shapiro, D. Dentcheva y A. Ruszczyński. \emph{Lectures on Stochastic Programming: Modeling and
Theory}. SIAM, 2009.

\bibitem{RU00} R. T. Rockafellar y S. Uryasev. Optimization of conditional value-at-risk. \emph{Journal of Risk},
2(3):21--41, 2000.

\bibitem{NJLS09} A. Nemirovski, A. Juditsky, G. Lan y A. Shapiro. Robust stochastic approximation approach to
stochastic programming. \emph{SIAM Journal on Optimization}, 19(4):1574--1609, 2009.

\bibitem{SSSS09} S. Shalev-Shwartz, O. Shamir, N. Srebro y K. Sridharan. Stochastic convex optimization.
\emph{Conference on Learning Theory (COLT)}, 2009.

\bibitem{Schwartz97} E. S. Schwartz. The stochastic behavior of commodity prices: Implications for valuation and
hedging. \emph{The Journal of Finance}, 52(3):923--973, 1997.

\bibitem{Rios24} D. Ríos, A. Pérez, J. Carabalí y L. Meneses. Poder de mercado y eventos climáticos adversos en un
mercado de electricidad hidro-dominado. \emph{Borradores de Economía} n.º 1266, Banco de la República (Colombia), 2024.

\bibitem{Valora24} Valora Analitik. «Niveles de embalses en Colombia llegan a mínimos de 20 años; piden medidas más
fuertes». Abril de 2024 (recoge la comunicación de XM, operador del sistema, al Gobierno nacional).
\url{https://www.valoraanalitik.com/niveles-de-embalses-en-colombia-llegan-a-minimos-de-20-anos/}

\bibitem{NOAA} NOAA Climate Prediction Center. Oceanic Niño Index (ONI): \emph{Cold \& warm episodes by season}.
\url{https://origin.cpc.ncep.noaa.gov/products/analysis_monitoring/ensostuff/ONI_v5.php}

\end{thebibliography}
\end{document}
